\documentclass[11pt,a4paper]{article}
\usepackage[a4paper,left=18mm,right=18mm,top=19mm,bottom=18mm,headheight=14pt,headsep=7mm,footskip=9mm]{geometry}
\usepackage{fontspec}
\IfFileExists{/usr/share/texmf/fonts/opentype/public/lm/lmroman10-regular.otf}{\setmainfont{lmroman10-regular.otf}[Path=/usr/share/texmf/fonts/opentype/public/lm/,BoldFont=lmroman10-bold.otf,ItalicFont=lmroman10-italic.otf,BoldItalicFont=lmroman10-bolditalic.otf]}{\setmainfont{Latin Modern Roman}}
\IfFileExists{/usr/share/fonts/truetype/dejavu/DejaVuSans.ttf}{\setsansfont{DejaVuSans.ttf}[Path=/usr/share/fonts/truetype/dejavu/,BoldFont=DejaVuSans-Bold.ttf,ItalicFont=DejaVuSans-Oblique.ttf]}{\setsansfont{Latin Modern Sans}}
\IfFileExists{/usr/share/texmf/fonts/opentype/public/lm/lmmono10-regular.otf}{\setmonofont{lmmono10-regular.otf}[Path=/usr/share/texmf/fonts/opentype/public/lm/,Scale=0.92,BoldFont=lmmonolt10-bold.otf]}{\setmonofont{Latin Modern Mono}[Scale=.92]}
\usepackage{amsmath,amssymb,booktabs,tabularx,longtable,array,enumitem,microtype,xcolor,graphicx,fancyhdr,titlesec,hyperref,xurl,lastpage,listings,tcolorbox,tikz}
\tcbuselibrary{skins,breakable}
\usetikzlibrary{arrows.meta,positioning}
\definecolor{Teal}{HTML}{1C5D66}\definecolor{Ink}{HTML}{20333C}\definecolor{Gold}{HTML}{B08043}\definecolor{Pale}{HTML}{F0F5F4}\definecolor{Line}{HTML}{CEDCDB}\definecolor{Muted}{HTML}{506770}
\color{Ink}\linespread{1.02}\setlength{\parindent}{0pt}\setlength{\parskip}{5pt}
\setlength{\emergencystretch}{2em}\clubpenalty=10000\widowpenalty=10000
\hypersetup{colorlinks=true,linkcolor=Teal,urlcolor=Teal,pdftitle={Five Nights at Gudo's - Project Study Guide},pdfauthor={Project documentation},bookmarksopen=true,bookmarksnumbered=false}
\urlstyle{tt}\Urlmuskip=0mu plus 1mu
\graphicspath{{figures/}{docs/defense/figures/}}
\setlist[itemize]{leftmargin=4mm,itemsep=2pt,topsep=3pt,parsep=0pt}
\setlist[enumerate]{leftmargin=5mm,itemsep=3pt,topsep=3pt,parsep=0pt}
\pagestyle{fancy}\fancyhf{}\fancyhead[L]{\sffamily\scriptsize FIVE NIGHTS AT GUDO'S}\fancyhead[R]{\sffamily\scriptsize PROJECT STUDY GUIDE}
\fancyfoot[L]{\sffamily\scriptsize Source snapshot: e9f0ff45}\fancyfoot[R]{\sffamily\scriptsize\thepage\,/\,\pageref*{LastPage}}
\renewcommand{\headrulewidth}{0.5pt}\renewcommand{\footrulewidth}{0pt}
\titleformat{\section}{\sffamily\bfseries\Large\color{Teal}}{}{0pt}{}[\vspace{2pt}{\color{Gold}\titlerule}]
\titleformat{\subsection}{\sffamily\bfseries\normalsize\color{Teal}}{}{0pt}{}
\titlespacing*{\section}{0pt}{0pt}{9pt}\titlespacing*{\subsection}{0pt}{7pt}{4pt}
\newcommand{\guidepage}[2]{\clearpage\phantomsection\label{#2}\pdfbookmark[1]{#1}{bm-#2}\section*{#1}}
\newcommand{\code}[1]{\nolinkurl{#1}}
\newcommand{\sources}[1]{\par\vspace{2pt}{\sffamily\fontsize{8}{10}\selectfont\color{Muted}\textbf{Open in source:} #1}\par}
\newcommand{\note}[2]{\begin{tcolorbox}[enhanced,colback=Pale,colframe=Line,boxrule=0.4pt,arc=1mm,left=7pt,right=7pt,top=5pt,bottom=5pt]\textbf{#1}\par#2\end{tcolorbox}}
\newcommand{\figurepair}[4]{\begin{center}\begin{minipage}{.485\linewidth}\includegraphics[width=\linewidth]{#1}\par{\sffamily\scriptsize #2}\end{minipage}\hfill\begin{minipage}{.485\linewidth}\includegraphics[width=\linewidth]{#3}\par{\sffamily\scriptsize #4}\end{minipage}\end{center}}
\newcommand{\filecard}[5]{\begin{tcolorbox}[enhanced,colback=white,colframe=Line,boxrule=.4pt,arc=1mm,left=8pt,right=8pt,top=6pt,bottom=6pt,before skip=6pt,after skip=6pt]\phantomsection\label{file:#1}\pdfbookmark[2]{#1.java}{bm-file-#1}{\sffamily\bfseries\color{Teal}#1.java}\hfill{\scriptsize #2}\par\small #3\par\textbf{Main path.} #4\par\textbf{Design choice.} #5\end{tcolorbox}}
\newcommand{\q}[2]{\par\vspace{4pt}{\sffamily\bfseries\color{Teal}#1}\par #2\par}
\lstset{language=Java,basicstyle=\ttfamily\small,keywordstyle=\color{Teal}\bfseries,commentstyle=\color{Muted},stringstyle=\color{Gold},backgroundcolor=\color{Pale},frame=single,rulecolor=\color{Line},breaklines=true,columns=fullflexible,keepspaces=true,showstringspaces=false,aboveskip=6pt,belowskip=6pt,framesep=6pt}
\newcolumntype{L}[1]{>{\raggedright\arraybackslash}p{#1}}\renewcommand{\arraystretch}{1.16}\setlength{\tabcolsep}{5pt}
\begin{document}
\pagenumbering{Roman}
\thispagestyle{empty}\pdfbookmark[0]{Quick topic index (unnumbered)}{front-index}
\section*{Quick topic index}
\textbf{Find the topic, then click its page number.} This front index is not part of the 42 numbered reviewer pages. References below and in source comments use the \textbf{printed page numbers}; the cover that follows is page 1.
\par\vspace{3pt}
\begin{minipage}[t]{.485\linewidth}\fontsize{9}{11.3}\selectfont
\begin{tabularx}{\linewidth}{@{}Xr@{}}\toprule\textbf{Topic / feature}&\textbf{Page}\\\midrule
\multicolumn{2}{@{}l}{\rule{0pt}{17pt}\textbf{\color{Teal}Start, navigate and demonstrate}}\\
Project overview / Java vocabulary&\hyperref[p:overview]{\pageref*{p:overview}}\\
Build and launch / find code&\hyperref[p:run]{\pageref*{p:run}}\\
Normal shift / controls walkthrough&\hyperref[p:normal]{\pageref*{p:normal}}\\
Settings / profiles / records&\hyperref[p:extras]{\pageref*{p:extras}}\\
Hidden game / five paper clicks&\hyperref[p:secret]{\pageref*{p:secret}}\\
\multicolumn{2}{@{}l}{\rule{0pt}{17pt}\textbf{\color{Teal}Rules and visible behavior}}\\
Clock / power / event deadlines&\hyperref[p:rules]{\pageref*{p:rules}}\\
Gudo route / difficulty / flashlight&\hyperref[p:gudo]{\pageref*{p:gudo}}\\
Input coordinates / panorama&\hyperref[p:render]{\pageref*{p:render}}\\
Pause / resume / abandonment&\hyperref[p:application]{\pageref*{p:application}}\\
Monitor / camera / accepted input&\hyperref[p:gamesession]{\pageref*{p:gamesession}}\\
HUD / camera feed / light drawing&\hyperref[p:gamepanel]{\pageref*{p:gamepanel}}\\
\multicolumn{2}{@{}l}{\rule{0pt}{17pt}\textbf{\color{Teal}Data and saving}}\\
SQLite tables / relationships&\hyperref[p:data]{\pageref*{p:data}}\\
Transactions / rollback / UUID retry&\hyperref[p:save]{\pageref*{p:save}}\\
Save folder / final run receipt&\hyperref[p:files-7]{\pageref*{p:files-7}}\\
Settings and shift-log text stores&\hyperref[p:files-12]{\pageref*{p:files-12}}\\
Records table / preference screens&\hyperref[p:files-6]{\pageref*{p:files-6}}\\
\multicolumn{2}{@{}l}{\rule{0pt}{17pt}\textbf{\color{Teal}Patterns and design decisions}}\\
New Game / startNight Facade&\hyperref[p:facade]{\pageref*{p:facade}}\\
AudioService Facade&\hyperref[p:facade]{\pageref*{p:facade}}\\
Renderer Decorator / effect order&\hyperref[p:decorator]{\pageref*{p:decorator}}\\
SecretLevelBuilder / Builder&\hyperref[p:builder]{\pageref*{p:builder}}\\
Template Method / inheritance / OOP&\hyperref[p:builder]{\pageref*{p:builder}}\\
Threads / EDT / audio resources&\hyperref[p:services]{\pageref*{p:services}}\\
\bottomrule\end{tabularx}\end{minipage}\hfill\begin{minipage}[t]{.485\linewidth}\fontsize{9}{11.3}\selectfont
\begin{tabularx}{\linewidth}{@{}Xr@{}}\toprule\textbf{Topic / feature}&\textbf{Page}\\\midrule
\multicolumn{2}{@{}l}{\rule{0pt}{17pt}\textbf{\color{Teal}Find an implementation}}\\
Application: navigation and saves&\hyperref[p:application]{\pageref*{p:application}}\\
GameSession: normal-night model&\hyperref[p:gamesession]{\pageref*{p:gamesession}}\\
GamePanel: input and rendering&\hyperref[p:gamepanel]{\pageref*{p:gamepanel}}\\
Main / MenuPanel / Navigation&\hyperref[p:files-1]{\pageref*{p:files-1}}\\
Extras / night selection / info&\hyperref[p:files-2]{\pageref*{p:files-2}}\\
ScenePanel / buttons / backdrop&\hyperref[p:files-3]{\pageref*{p:files-3}}\\
Theme / loading / VHS overlay&\hyperref[p:files-4]{\pageref*{p:files-4}}\\
Pause / result / jumpscare&\hyperref[p:files-5]{\pageref*{p:files-5}}\\
Settings / RecordsPanel&\hyperref[p:files-6]{\pageref*{p:files-6}}\\
Paths / Progress / RunRecord&\hyperref[p:files-7]{\pageref*{p:files-7}}\\
Scene / Outcome / CameraMotion&\hyperref[p:files-8]{\pageref*{p:files-8}}\\
Assets / panorama / camera renderer&\hyperref[p:files-9]{\pageref*{p:files-9}}\\
Renderer interface and wrappers&\hyperref[p:files-10]{\pageref*{p:files-10}}\\
Darkness / AudioService / Sound&\hyperref[p:files-11]{\pageref*{p:files-11}}\\
Chiptune / settings / shift log&\hyperref[p:files-12]{\pageref*{p:files-12}}\\
Database / SecretLevel / Builder&\hyperref[p:files-13]{\pageref*{p:files-13}}\\
GudoEncounter / secret runtime&\hyperref[p:files-14]{\pageref*{p:files-14}}\\
\multicolumn{2}{@{}l}{\rule{0pt}{17pt}\textbf{\color{Teal}Questions and presentation}}\\
Questions 1--10: overview / flow&\hyperref[p:qa-1]{\pageref*{p:qa-1}}\\
Questions 11--20: patterns / Java&\hyperref[p:qa-2]{\pageref*{p:qa-2}}\\
Questions 21--30: rules / timing&\hyperref[p:qa-3]{\pageref*{p:qa-3}}\\
Questions 31--40: Gudo / secret&\hyperref[p:qa-4]{\pageref*{p:qa-4}}\\
Questions 41--50: collision / data&\hyperref[p:qa-5]{\pageref*{p:qa-5}}\\
Questions 51--60: failures / limits&\hyperref[p:qa-6]{\pageref*{p:qa-6}}\\
Five-minute rehearsal / snags&\hyperref[p:rehearse]{\pageref*{p:rehearse}}\\
Limitations / source evidence&\hyperref[p:limits]{\pageref*{p:limits}}\\
Alphabetical index: all 45 Java files&\hyperref[p:index]{\pageref*{p:index}}\\
\bottomrule\end{tabularx}\end{minipage}
\par\vspace{7pt}
\note{If you are asked to explain a feature}{Open the topic page, then find its Java file using page 42. Follow this order: \textbf{what it does $\rightarrow$ where it lives $\rightarrow$ how it works $\rightarrow$ why the choice helps}. For New Game, start with page 13, then Application on page 17.}
\clearpage\pagenumbering{arabic}
\thispagestyle{empty}\pdfbookmark[0]{Cover}{cover}
\vspace*{7mm}
{\sffamily\small\color{Gold}\textbf{JAVA DESKTOP PROJECT\quad /\quad STUDY EDITION}}\par\vspace{7mm}
{\sffamily\bfseries\fontsize{34}{40}\selectfont FIVE NIGHTS\\AT GUDO'S}\par\vspace{5mm}
{\sffamily\Large\color{Teal}Architecture, source code and demonstration guide}\par\vspace{4mm}
A compact reference for understanding the implementation, finding each feature, explaining the important decisions and preparing for project questions.
\vspace{6mm}
\par\noindent\includegraphics[width=\linewidth]{office.jpg}\par
\vspace{5mm}
\par\noindent\begin{tabularx}{\linewidth}{XXX}\toprule
\textbf{45 application Java files}&\textbf{Runtime features}&\textbf{One coherent walkthrough}\\
Rules, screens, rendering, audio and persistence&Gameplay, cameras, sound and local data&Click paths, worked examples and source links\\\bottomrule\end{tabularx}\par
\vfill
\note{How this edition stays concise}{Every Java file is covered. Repeated getter explanations and full per-method narration are kept in the source comments; the guide concentrates on responsibilities, important methods, execution order, reasons and demonstrations.}
{\sffamily\small 15 November 2026\quad\textbullet\quad Checked source: \code{38bad499}}\par
\guidepage{Reading map}{p:map}

\subsection*{Choose a route}
\par\noindent\begin{tabularx}{\linewidth}{L{.27\linewidth}XL{.13\linewidth}}\toprule
\textbf{Need}&\textbf{Read}&\textbf{Page}\\\midrule
Start from zero&Project vocabulary and ownership&\pageref{p:overview}\\
Run and navigate&Build, locate code, normal/secret demonstrations&\pageref{p:run}--\pageref{p:secret}\\
Explain the rules&Clock/power, Gudo, input and projection&\pageref{p:rules}--\pageref{p:render}\\
Explain saves&Schema, transaction, text files and retry&\pageref{p:data}--\pageref{p:save}\\
Explain patterns&Both facades, Decorator, Builder and OOP&\pageref{p:facade}--\pageref{p:builder}\\
Trace major classes&Application, GameSession and GamePanel&\pageref{p:application}--\pageref{p:gamepanel}\\
Find a Java file&All application source-file cards&\pageref{p:files-1}--\pageref{p:files-14}\\
Practice questions&Sixty short, code-grounded answers&\pageref{p:qa-1}--\pageref{p:qa-6}\\
Rehearse&Presentation sequence and limits&\pageref{p:rehearse}--\pageref{p:limits}\\
Locate quickly&Alphabetical file and page index&\pageref{p:index}\\\bottomrule
\end{tabularx}\par
\subsection*{A useful answer structure}
\textbf{What it does $\rightarrow$ where it lives $\rightarrow$ how the code works $\rightarrow$ why the arrangement helps.}
For example: ``The flashlight protects the office. GameSession and GudoEncounter enforce the rule. Correct-facing light builds exposure until a threshold; interruption reduces it. Keeping this outside drawing code makes the rule testable.''
\subsection*{Use the clickable references}
A source link opens the exact annotated revision. In an IDE, expand \code{src/main/java/com/gudos}, open the named file and search for the method declaration with Ctrl+F. If line numbers change later, the exact filename and method name still locate the feature.
\note{The three facts that prevent confusing answers}{New Game preserves history. Continue starts a fresh attempt rather than restoring a checkpoint. A hidden-minigame win does not unlock the next normal night.}
\subsection*{Scope of explanations}
Rationales describe benefits and tradeoffs visible in the implementation, not undocumented claims about the original author's intentions. Screenshot examples are actual Swing renders with controlled demonstration state; they do not certify native fullscreen, focus handling or speaker output on every computer.

\clearpage\pdfbookmark[0]{Foundations and navigation}{group-p:overview}\phantomsection\label{p:overview}\pdfbookmark[1]{The whole project in one page}{bm-p:overview}\section*{The whole project in one page}

\textbf{Purpose.} An offline Java survival game set in a carinderia. Last six active minutes, displayed as 12 AM--6 AM, manage power and repel one enemy. Five normal nights unlock progressively; a hidden platformer is a separate interaction.
\par\noindent\begin{tabularx}{\linewidth}{L{.25\linewidth}X}\toprule\textbf{Owner}&\textbf{Responsibility}\\\midrule
\code{Main}&Creates the loading window and prepares dependencies.\\
\code{Application}&Coordinates screens, attempts, audio, storage and orderly exit.\\
\code{GameSession / GudoEncounter}&Normal-night rules: time, power, routes, directional defense and final outcome.\\
\code{GamePanel / SecretPanel}&Translate input and draw the current model.\\
\code{SecretGame / SecretLevel}&Changing platformer state versus its configuration.\\
\code{Database / text stores}&Local progress/history, preferences and readable reports.\\
\code{AssetManager / renderers / AudioService}&Cached resources, visual composition and sound playback.\\\bottomrule\end{tabularx}\par
\subsection*{Essential Java vocabulary}
\begin{description}[leftmargin=19mm,style=nextline,itemsep=2pt]
\item[Class / object] A type and one instance of that type: \code{new GameSession(1)} creates an attempt.
\item[Method / field] An operation versus retained data: \code{toggleMonitor()} changes accepted monitor state; \code{power} stores a number.
\item[Interface] A small callable contract: panels depend on \code{Navigation}; renderers share \code{SceneRenderer}.
\item[Callback] Behavior supplied for later use: \code{() -> navigation.startNight(1)} runs when the button is activated.
\item[Enum / record] Named finite values versus concise data carriers: \code{Outcome.CAUGHT}; \code{RunRecord}.
\item[EDT / executor] Swing's UI thread versus a worker queue. UI changes return to EDT; IO and sound use workers.
\item[Transaction / retry] Related database changes commit together; an idempotent retry repeats a request without duplicating its intended result.
\end{description}
\note{Keep these values separate}{\textbf{Settings} are preferences. \textbf{Progress} is the loaded unlock/count summary. \textbf{GameSession} is one live attempt. \textbf{RunRecord} is its immutable final receipt, not a resumable live save.}
\sources{\href{https://github.com/itsnold/FiveNightsAtGudos/blob/e9f0ff45350572e7ba5fd0ff6185ab1251514e2f/src/main/java/com/gudos/Main.java\#L31}{\nolinkurl{Main.java}:31}; \href{https://github.com/itsnold/FiveNightsAtGudos/blob/e9f0ff45350572e7ba5fd0ff6185ab1251514e2f/src/main/java/com/gudos/Application.java\#L51}{\nolinkurl{Application.java}:51}; \href{https://github.com/itsnold/FiveNightsAtGudos/blob/e9f0ff45350572e7ba5fd0ff6185ab1251514e2f/src/main/java/com/gudos/game/GameSession.java\#L28}{\nolinkurl{game/GameSession.java}:28}; \href{https://github.com/itsnold/FiveNightsAtGudos/blob/e9f0ff45350572e7ba5fd0ff6185ab1251514e2f/src/main/java/com/gudos/game/RunRecord.java\#L23}{\nolinkurl{game/RunRecord.java}:23}}
\guidepage{Build, launch and find the implementation}{p:run}

\subsection*{Windows preparation}
\begin{enumerate}
\item Open the repository root containing \code{pom.xml}, \code{run.cmd}, \code{mvnw.cmd} and \code{src}. Install a JDK 17 or newer; building needs compiler tooling.
\item Build the application before the demonstration. First build needs internet for Maven/dependencies.
\end{enumerate}
\begin{lstlisting}[language={}]
.\mvnw.cmd clean package
\end{lstlisting}
\begin{enumerate}[resume]
\item Double-click \code{run.cmd} for ordinary launch. It chooses a JDK, packages the application and runs the game. The command above performs the same packaging explicitly.
\item For an isolated demo, launch the distributable with its own data directory:
\end{enumerate}
\begin{lstlisting}[language={}]
java --enable-native-access=ALL-UNNAMED `
  "-Dgudo.dataDir=target/defense-demo-data" `
  -jar target/five-nights-at-gudos.jar
\end{lstlisting}
The backticks continue a command in PowerShell. On macOS/Linux build with \code{sh ./mvnw clean package}; macOS also has \code{run.command}. The shaded \code{five-nights-at-gudos.jar} bundles dependencies/resources and runs offline after building.
\subsection*{Find a feature, step by step}
\begin{enumerate}
\item Open the entire root as a Maven project. Expand \code{src/main/java/com/gudos}.
\item Choose \code{game} for rules, \code{ui} for screens, \code{persistence} for saves, \code{assets} for images/projection or \code{audio} for playback.
\item To explain a button, open its panel constructor and locate the callback passed to \code{GameButton}. Follow that callback into \code{Application} or the model.
\item Search the exact method name. Read its guard first, then changed state, then called operations. Use Go to Definition if available.
\item Follow the call into the model and identify its guard, state changes and next operation. Those implementation details are the evidence for the behavior.
\end{enumerate}
\subsection*{Explain the build accurately}
\code{pom.xml} targets Java 17. SQLite JDBC supplies embedded storage and the selected SLF4J backend handles logging integration. The compiler builds application classes; Shade produces the distributable with \code{com.gudos.Main} as its manifest entry point. Swing/Java2D provide the UI and drawing stack; no browser or separate game engine is involved.
\sources{\href{https://github.com/itsnold/FiveNightsAtGudos/blob/e9f0ff45350572e7ba5fd0ff6185ab1251514e2f/src/main/java/com/gudos/Main.java\#L43}{\nolinkurl{Main.java}:43}; \href{https://github.com/itsnold/FiveNightsAtGudos/blob/e9f0ff45350572e7ba5fd0ff6185ab1251514e2f/src/main/java/com/gudos/Main.java\#L54}{\nolinkurl{Main.java}:54}}
\guidepage{Normal shift: the exact demonstration path}{p:normal}

\figurepair{menu.jpg}{Title: New Game, Continue and Extras}{office.jpg}{Office: clock, power, light and monitor tab}
\begin{enumerate}
\item \textbf{New Game.} Click it, or use Up/Down and Enter. It calls \code{navigation.startNight(1)}. Application checks guards, constructs a fresh model/view, wires callbacks, switches screen and requests audio. History remains.
\item \textbf{Look around.} Hold A/D or arrows, or move the pointer near either horizontal edge above the bottom controls. Office pan eases smoothly instead of snapping.
\item \textbf{CCTV.} Press Space/Tab or the bottom chevron tab. Press 1--6 or click CAM buttons. Each feed sweeps automatically using the active shift clock, not the mouse.
\item \textbf{Light.} Lower the monitor; hold F or the light button, then release. Raised monitor blocks light in the model. Holding F while facing the wrong door still costs power.
\item \textbf{First defense.} On Night 1: CAM04 at 18 s, CAM03 at 29 s, CAM02 at 40 s, left office at 51 s. Lower the monitor, face left and hold F for at least 1.6 effective seconds. Release after he disappears.
\item \textbf{Pause.} Escape opens PausePanel. Escape/Resume restores the same in-memory attempt. Clock, drain, enemy and camera sweep stop while paused.
\item \textbf{Leave.} Escape $\rightarrow$ RETURN TO MENU $\rightarrow$ confirm Yes. The attempt becomes ABANDONED, saves, then returns to title. It does not unlock a night.
\end{enumerate}
\note{Continue and replay}{Continue enables after any recorded attempt and starts \code{progress.nextNight()} fresh. Extras $\rightarrow$ Select Night offers unlocked nights; Application rechecks eligibility. Replaying an earlier night cannot reduce saved unlocks.}
\textbf{Ending demonstration.} Normal victory needs six active minutes. Ignoring the first Night 1 encounter naturally catches at 60 s, then a 2.1 s scare plays while saving starts. Power failure is separate; there is no runtime ``force power failure'' button.
\sources{\href{https://github.com/itsnold/FiveNightsAtGudos/blob/e9f0ff45350572e7ba5fd0ff6185ab1251514e2f/src/main/java/com/gudos/ui/MenuPanel.java\#L86}{\nolinkurl{ui/MenuPanel.java}:86}; \href{https://github.com/itsnold/FiveNightsAtGudos/blob/e9f0ff45350572e7ba5fd0ff6185ab1251514e2f/src/main/java/com/gudos/Application.java\#L306}{\nolinkurl{Application.java}:306}; \href{https://github.com/itsnold/FiveNightsAtGudos/blob/e9f0ff45350572e7ba5fd0ff6185ab1251514e2f/src/main/java/com/gudos/ui/GamePanel.java\#L266}{\nolinkurl{ui/GamePanel.java}:266}; \href{https://github.com/itsnold/FiveNightsAtGudos/blob/e9f0ff45350572e7ba5fd0ff6185ab1251514e2f/src/main/java/com/gudos/Application.java\#L436}{\nolinkurl{Application.java}:436}}
\guidepage{Preferences, records and profile navigation}{p:extras}

\figurepair{settings.jpg}{Extras $\rightarrow$ Settings}{records.jpg}{Extras $\rightarrow$ Shift Records}
\subsection*{Apply settings and select a local profile}
\begin{enumerate}
\item Title $\rightarrow$ Extras $\rightarrow$ Settings. Edit player name, volume, reduced effects or fullscreen.
\item Click APPLY SETTINGS and wait for title. Application loads/creates the named profile and writes preferences on IO; success replaces loaded settings/progress and updates volume/display.
\item Reopen Settings or restart to confirm persistence. Back without Apply discards form edits. F11 separately persists/toggles fullscreen.
\item A different name selects separate local records. Names are not authenticated accounts. The model removes control characters, trims/caps names and clamps volume.
\end{enumerate}
\subsection*{Locate the actual save files}
\begin{enumerate}
\item Settings $\rightarrow$ COPY DATA FOLDER. Paste into File Explorer's address bar. The explicit demo-directory override wins over ordinary locations.
\item Open \code{settings.txt} for preferences and \code{shift-log.txt} for readable receipts. Inspect a copy of \code{gudo.db} in a SQLite viewer if needed.
\item Title $\rightarrow$ Extras $\rightarrow$ Shift Records. Application loads rows off EDT; RecordsPanel displays an already-loaded list. The table is read-only.
\end{enumerate}
\note{Totals versus rows}{The screen shows the latest 200 attempts, but older rows remain stored and totals count them. Only normal COMPLETED outcomes increase completed counts/progression. Names use SQLite NOCASE behavior, not full Unicode account identity.}
\textbf{Other Extras routes.} Select Night gates buttons from Progress; How to Play and Credits are static InfoPanel variants. Back/Escape on these secondary pages returns to the title, not to a previous-page history stack.
\sources{\href{https://github.com/itsnold/FiveNightsAtGudos/blob/e9f0ff45350572e7ba5fd0ff6185ab1251514e2f/src/main/java/com/gudos/Application.java\#L274}{\nolinkurl{Application.java}:274}; \href{https://github.com/itsnold/FiveNightsAtGudos/blob/e9f0ff45350572e7ba5fd0ff6185ab1251514e2f/src/main/java/com/gudos/Application.java\#L529}{\nolinkurl{Application.java}:529}; \href{https://github.com/itsnold/FiveNightsAtGudos/blob/e9f0ff45350572e7ba5fd0ff6185ab1251514e2f/src/main/java/com/gudos/ui/SettingsPanel.java\#L192}{\nolinkurl{ui/SettingsPanel.java}:192}; \href{https://github.com/itsnold/FiveNightsAtGudos/blob/e9f0ff45350572e7ba5fd0ff6185ab1251514e2f/src/main/java/com/gudos/persistence/Database.java\#L283}{\nolinkurl{persistence/Database.java}:283}}
\guidepage{Hidden game: trigger, rules and restart}{p:secret}

\figurepair{camera04.jpg}{CAM04 stockroom: click the small left-side paper}{secret.jpg}{The scrolling hidden minigame}
\begin{enumerate}
\item Start Night 1--4, raise monitor and choose CAM04. Left-click the actual paper five valid times. Follow its position as the feed sweeps; camera-map clicks are not the trigger.
\item The code reverse-maps the click into source-image pixels. Valid bounds are x=28.5--32.5\% and y=24--40\% of source size. Invalid clicks neither increment nor reset valid count. The fifth valid click runs the callback. Application rejects Night 5.
\item Application releases normal controls, pauses the normal session, stops audio and constructs the hidden model/panel. Construction failure restores the existing normal shift.
\item A/D or arrows move, Space jumps and E picks up/drops/delivers. Escape pauses this minigame; focus loss also pauses and clears held keys.
\item Explore before pickup freely. The first successful pickup sets \code{started=true}; 60 s then count down even after dropping or hiding. Re-pickup never resets time.
\item Witnesses detect the bundle center (carried or dropped) in their triangular beams unless a cover rectangle contains it. The player's unladen body alone does not cause detection. The view draws the same beam polygon the model checks.
\item Reach the right-hand stockroom while carrying and press E. Delivery is checked before dropping. Timeout or beam detection loses.
\item Both outcomes stop sound and freeze with fixed signal lines for three unpaused seconds. Application saves the secret outcome and waits for both stores, then starts a \emph{fresh attempt of the same normal night}.
\end{enumerate}
\note{Important distinction}{MINIGAME\_COMPLETED is not normal COMPLETED. It does not unlock the next night and does not permanently prevent replay. The misleadingly named \code{nextNightAfterSecret} stores the same night to restart.}
\sources{\href{https://github.com/itsnold/FiveNightsAtGudos/blob/e9f0ff45350572e7ba5fd0ff6185ab1251514e2f/src/main/java/com/gudos/ui/GamePanel.java\#L183}{\nolinkurl{ui/GamePanel.java}:183}; \href{https://github.com/itsnold/FiveNightsAtGudos/blob/e9f0ff45350572e7ba5fd0ff6185ab1251514e2f/src/main/java/com/gudos/Application.java\#L348}{\nolinkurl{Application.java}:348}; \href{https://github.com/itsnold/FiveNightsAtGudos/blob/e9f0ff45350572e7ba5fd0ff6185ab1251514e2f/src/main/java/com/gudos/game/minigame/SecretGame.java\#L92}{\nolinkurl{game/minigame/SecretGame.java}:92}; \href{https://github.com/itsnold/FiveNightsAtGudos/blob/e9f0ff45350572e7ba5fd0ff6185ab1251514e2f/src/main/java/com/gudos/Application.java\#L396}{\nolinkurl{Application.java}:396}}
\clearpage\pdfbookmark[0]{Rules and persistence}{group-p:rules}\phantomsection\label{p:rules}\pdfbookmark[1]{Clock, power and event-boundary updates}{bm-p:rules}\section*{Clock, power and event-boundary updates}

\subsection*{Six active minutes, six fictional hours}
A normal duration is $D=360$ seconds. Progress is $\min(1,t/D)$. The HUD maps it to 12 AM--6 AM; preciseClock adds fictional minutes. RunRecord duration instead formats actual active seconds. Wall-clock Instants can be far apart after a long pause without increasing active duration.
\subsection*{Power is a rate, not a bar count}
With monitor flag $M$ and light flag $L$:
\[
 d=0.085+0.22M+0.20L,\qquad P_{new}=\max(0,P-d\Delta t).
\]
The model excludes simultaneous $M=L=1$. Idle costs 0.085, monitor 0.305 and light 0.285 percentage points per active second. Ten idle plus ten monitor seconds leave $100-0.85-3.05=96.10\%$.
\par\noindent\begin{tabularx}{\linewidth}{L{.23\linewidth}XL{.23\linewidth}}\toprule
\textbf{Example}&\textbf{Calculation}&\textbf{Meaning}\\\midrule
Idle full shift&$100-0.085(360)=69.4$&Defense adds extra light cost.\\
Always-monitor budget&$100/0.305\approx327.87$ s&Gudo may catch earlier.\\
Usage display&$1+2M+2L$&Normally 1 or 3 bars, not drain.\\\bottomrule\end{tabularx}\par
\subsection*{Why a frame is split into smaller events}
\begin{enumerate}
\item Reject a negative/nonfinite delta; ignore paused or finished sessions.
\item Compute morning time remaining, power-depletion time and Gudo's next boundary. Effective light also requires the correct approach direction.
\item Advance by the minimum of those deadlines and remaining frame time. Subtract drain for that exact step.
\item Test morning first, power failure second; otherwise advance the encounter and check caught. Repeat for the rest of the frame.
\end{enumerate}
This prevents a large supplied delta from skipping an arrival, attack or repel. On a simultaneous deadline morning wins, then power failure, then attack. \code{finish} freezes one immutable RunRecord and disables equipment; later finish requests cannot change it.
\note{UI timer versus model accuracy}{Application requests ticks around every 33 ms, measures nanoTime and caps one delayed delta at 0.25 s. Severe stalls can make simulation lag real wall time. Within each supplied delta, the model still handles exact boundaries.}
\sources{\href{https://github.com/itsnold/FiveNightsAtGudos/blob/e9f0ff45350572e7ba5fd0ff6185ab1251514e2f/src/main/java/com/gudos/game/GameSession.java\#L80}{\nolinkurl{game/GameSession.java}:80}; \href{https://github.com/itsnold/FiveNightsAtGudos/blob/e9f0ff45350572e7ba5fd0ff6185ab1251514e2f/src/main/java/com/gudos/game/GameSession.java\#L215}{\nolinkurl{game/GameSession.java}:215}; \href{https://github.com/itsnold/FiveNightsAtGudos/blob/e9f0ff45350572e7ba5fd0ff6185ab1251514e2f/src/main/java/com/gudos/game/GameSession.java\#L111}{\nolinkurl{game/GameSession.java}:111}}
\guidepage{Gudo: routes, difficulty and effective light}{p:gudo}

\par\noindent\begin{tabularx}{\linewidth}{L{.15\linewidth}XL{.22\linewidth}}\toprule
\textbf{Approach}&\textbf{Route then office}&\textbf{Facing condition}\\\midrule
LEFT&04 stockroom $\rightarrow$03 kitchen $\rightarrow$02 counter&$p\leq0.30$; target 0.12\\
RIGHT&06 alley $\rightarrow$05 corridor&$p\geq0.70$; target 0.88\\
FRONT&01 front entrance&$|p-0.50|\leq0.22$\\\bottomrule\end{tabularx}\par
Pan $p$ is normalized: 0=left, 0.5=center, 1=right. Only the current approach is checked; small edge overlaps do not create multiple simultaneous enemies. CAM01 currently means front entrance, despite some retained dining filenames/older prose.
\subsection*{State transition, in order}
Waiting uses routeIndex$=-1$ and location=null. When timer expires, advance to the next camera; after the route, enter office with exposure 0 and an attack countdown. At office, correct light adds exposure. Releasing/looking away removes exposure at 0.65 per second, floored at zero. Enough exposure repels; otherwise expired timer catches.
After repel: clear occupancy, reset exposure/routeIndex, increment repelled and cycle LEFT $\rightarrow$ RIGHT $\rightarrow$ FRONT. Location is singular, so the previous room automatically becomes clean/custom again.
\par\noindent\begin{tabularx}{\linewidth}{ccccc}\toprule
\textbf{Night}&\textbf{First wait}&\textbf{Initial route step}&\textbf{Attack window}&\textbf{Light needed}\\\midrule
1&18.0 s&11.0 s&9.0 s&1.6 s\\2&16.0&9.7&8.2&1.8\\3&14.0&8.4&7.4&2.0\\4&12.0&7.1&6.6&2.2\\5&10.0&5.8&5.8&2.4\\\bottomrule\end{tabularx}\par
\[
\begin{aligned}
T_{flash}&=1.6+0.2(n-1), &T_{attack}&=9-0.8(n-1),\\
T_{step}&=\max(4,11-1.3(n-1)-\min(3,0.3r)),\\
T_{wait}&=\max(5,13-(n-1)-\min(4,r)).
\end{aligned}
\]
Here $n$ is night and $r$ is repelled count. First-wait formula is $18-2(n-1)$; later waits use the separate formula above.
\note{Worked interrupted defense}{Night 1: 0.8 s effective exposure gives 50\% fade. Release for 0.4 s loses $0.4(0.65)=0.26$, leaving 0.54 s. Another 1.06 effective seconds is needed, while the attack countdown continues.}
\textbf{First Night 1 timeline:}18 s CAM04;29 s CAM03;40 s CAM02;51 s left office; earliest immediate repel 52.6 s; unattended caught 60 s.
\sources{\href{https://github.com/itsnold/FiveNightsAtGudos/blob/e9f0ff45350572e7ba5fd0ff6185ab1251514e2f/src/main/java/com/gudos/game/GudoEncounter.java\#L114}{\nolinkurl{game/GudoEncounter.java}:114}; \href{https://github.com/itsnold/FiveNightsAtGudos/blob/e9f0ff45350572e7ba5fd0ff6185ab1251514e2f/src/main/java/com/gudos/game/GudoEncounter.java\#L72}{\nolinkurl{game/GudoEncounter.java}:72}; \href{https://github.com/itsnold/FiveNightsAtGudos/blob/e9f0ff45350572e7ba5fd0ff6185ab1251514e2f/src/main/java/com/gudos/game/GameSession.java\#L191}{\nolinkurl{game/GameSession.java}:191}}
\guidepage{Input, scaling and image projection}{p:render}

\subsection*{Two coordinate transforms}
ScenePanel defines a logical $1280\times720$ scene. For physical width $w$, height $h$:
\[
 s=\min(w/1280,h/720),\qquad
 x_s=\frac{x-o_x}{s},\quad y_s=\frac{y-o_y}{s}.
\]
Offsets center the scaled scene; doLayout applies the same transform to real Swing controls. Storing original bounds/font avoids cumulative resize error. scenePoint reverses the transform for mouse input.
\subsection*{CCTV and office use different movement}
Office pan is requested by keys/pointer edges and eased with a factor $1-e^{-9\Delta t}$. The panel then sends it to GameSession for directional checks. CCTV instead uses CameraMotion:
\[
 p=0.5-0.46\cos(2\pi\phi),\quad
 \phi=t/(22+2.3c)+0.137c.
\]
Camera number $c$ changes phase/period; $p\in[0.04,0.96]$. Session time makes unseen feeds continue sweeping and all sweeps freeze on pause.
\subsection*{A panorama effect, not a 3D room}
PanoramaRenderer uses atan for horizontal rays and cos for vertical sampling near edges. It writes a reused $854\times480$ image scaled to the logical viewport. Source pixels/lookups are cached; pan is quantized to2049 positions. This differs from Theme.cover, which simply fills/crops an image rectangle. HUD text never goes through room projection.
\textbf{Buffer ownership.} render returns the same mutable output. Clean and occupied room images therefore need separate projection instances to coexist during a fade. Occupied alpha $=1-$fade is blended over clean.
\subsection*{Why the paper click stays attached to the artwork}
Convert event pixels $\rightarrow$logical scene $\rightarrow$projection pixel. Recompute CAM04 at current elapsed sweep, call sourcePoint and test the original-image fractional hitbox. A fixed window rectangle would drift during sweep/resizing.
\note{Reduced effects preserve rules}{Optional flicker, noise, blinking, intro and scare motion are reduced/bypassed. A disrupted feed still shows a static SIGNAL LOST screen for its two active seconds, preserving information hiding. Gameplay timing and power are unchanged.}
\sources{\href{https://github.com/itsnold/FiveNightsAtGudos/blob/e9f0ff45350572e7ba5fd0ff6185ab1251514e2f/src/main/java/com/gudos/ui/ScenePanel.java\#L89}{\nolinkurl{ui/ScenePanel.java}:89}; \href{https://github.com/itsnold/FiveNightsAtGudos/blob/e9f0ff45350572e7ba5fd0ff6185ab1251514e2f/src/main/java/com/gudos/assets/PanoramaRenderer.java\#L68}{\nolinkurl{assets/PanoramaRenderer.java}:68}; \href{https://github.com/itsnold/FiveNightsAtGudos/blob/e9f0ff45350572e7ba5fd0ff6185ab1251514e2f/src/main/java/com/gudos/ui/GamePanel.java\#L183}{\nolinkurl{ui/GamePanel.java}:183}}
\guidepage{Three save files and the relational schema}{p:data}

\par\noindent\begin{tabularx}{\linewidth}{L{.23\linewidth}X}\toprule\textbf{File}&\textbf{Actual responsibility}\\\midrule
\code{settings.txt}&UTF-8 preference keys: playerName, volume, reducedEffects, fullscreen.\\
\code{gudo.db}&Authoritative local profiles, unlock summary and ended attempt history.\\
\code{shift-log.txt}&Readable per-attempt report; not a checkpoint or progress replacement.\\\bottomrule\end{tabularx}\par
\textbf{Directory priority.} Explicit \code{gudo.dataDir}; else LOCALAPPDATA/FiveNightsAtGudos when that environment value exists; else user.home/.five-nights-at-gudos. Files are outside the read-only JAR. SQLite WAL/SHM sidecars are normal in WAL mode.
\subsection*{Tables and relationships}
\par\noindent\begin{tabularx}{\linewidth}{L{.19\linewidth}X}\toprule
\textbf{Table}&\textbf{Keys, data and constraints}\\\midrule
profiles&id INTEGER primary key; name required, UNIQUE COLLATE NOCASE; created\_at required.\\
progress&profile\_id primary/foreign key; next\_night 1--5 default 1; highest\_completed 0--5 default 0; updated\_at required.\\
night\_runs&id TEXT UUID primary key; profile\_id foreign key; night 1--5; elapsed\_seconds$\geq0$; power\_remaining 0--100; six allowed outcome names; start/end Instant text.\\\bottomrule\end{tabularx}\par
One initialized profile has one progress row and many runs. Index \code{runs_by_profile(profile_id,ended_at DESC)} supports the recent-history query. \code{loadRuns} returns newest 200, while COUNT queries include older retained rows.
\subsection*{Constraint vocabulary}
\textbf{PK}: unique row identity. \textbf{FK}: points to another table's row. \textbf{NOT NULL}: value required. \textbf{UNIQUE}: duplicates rejected under that comparison. \textbf{CHECK}: value must satisfy a condition. \textbf{Index}: lookup/order structure, not a new copy of game logic.
\subsection*{Connections and safe values}
Each operation uses a short-lived connection, enabling foreign keys and a 5000 ms busy timeout. PreparedStatement binds player values separately from SQL syntax, including quoted names. It does not provide encryption/authentication. An internally formatted migration table name is not user input.
\note{Local identity, not an account system}{Changing playerName selects/creates a local profile. SQLite's built-in NOCASE is not complete Unicode case folding. There is no password table, online authentication or synchronization here.}
\sources{\href{https://github.com/itsnold/FiveNightsAtGudos/blob/e9f0ff45350572e7ba5fd0ff6185ab1251514e2f/src/main/java/com/gudos/persistence/AppPaths.java\#L40}{\nolinkurl{persistence/AppPaths.java}:40}; \href{https://github.com/itsnold/FiveNightsAtGudos/blob/e9f0ff45350572e7ba5fd0ff6185ab1251514e2f/src/main/java/com/gudos/persistence/Database.java\#L133}{\nolinkurl{persistence/Database.java}:133}; \href{https://github.com/itsnold/FiveNightsAtGudos/blob/e9f0ff45350572e7ba5fd0ff6185ab1251514e2f/src/main/java/com/gudos/persistence/Database.java\#L154}{\nolinkurl{persistence/Database.java}:154}; \href{https://github.com/itsnold/FiveNightsAtGudos/blob/e9f0ff45350572e7ba5fd0ff6185ab1251514e2f/src/main/java/com/gudos/persistence/Database.java\#L203}{\nolinkurl{persistence/Database.java}:203}}
\guidepage{Save, retry, migration and preference writes}{p:save}

\subsection*{An ended attempt, from model to disk}
\begin{enumerate}
\item GameSession finishes once, creating its stable UUID RunRecord. Application retains it and the profile captured when the shift began.
\item ResultPanel disables conflicting navigation. The IO worker selects/creates profile/progress, then begins the attempt/unlock transaction.
\item INSERT night\_runs uses ON CONFLICT(id) DO NOTHING. If newly inserted \emph{and} outcome=COMPLETED, update unlock fields using MAX and cap next night at 5.
\item Commit those two database changes together, or rollback on SQLException. Close connection and reload Progress. Earlier profile setup is outside this transaction.
\item Append the readable log line after checking its run=UUID marker. Return success to EDT only when both operations finish.
\item Enable next/menu or honor queued exit/menu/same-night restart. Failure retains the pending result and exposes Retry Save.
\end{enumerate}
\note{A partial save is possible}{SQLite may commit before text append fails. Retry uses the same UUID: the database conflict inserts nothing, and the log marker prevents the ordinary duplicate report. This is retry coordination across separate stores, not one atomic cross-file transaction. A crash during plain append may leave a partial line.}
\subsection*{Replay and result invariants}
Already next\_night4 plus a Night 1 replay stays4 because MAX(4,2)=4. Completed Night 5 stays capped at 5. CAUGHT, POWER\_OUT, ABANDONED and both secret outcomes record history but do not unlock a normal night. Normal close can save ABANDONED after confirmation; force-kill cannot guarantee an unfinished attempt is saved.
\subsection*{Schema and text operations}
Database rejects schema versions above5. Supported legacy outcome-constrained tables are rebuilt/copied inside migration transactions; old rows and progress are preserved. Legacy secret flags remain but no longer gate replay. The entire constructor is not one all-encompassing transaction.
SettingsStore writes all keys to a temporary neighbor, then atomically replaces settings.txt where supported; otherwise regular replacement is used and temporary residue is cleaned. Loading uses parse fallbacks plus Settings normalization; IO failures still propagate. ShiftLogStore appends reports and scans existing UUID markers; it is not multi-process crash-proof journaling.
\sources{\href{https://github.com/itsnold/FiveNightsAtGudos/blob/e9f0ff45350572e7ba5fd0ff6185ab1251514e2f/src/main/java/com/gudos/Application.java\#L486}{\nolinkurl{Application.java}:486}; \href{https://github.com/itsnold/FiveNightsAtGudos/blob/e9f0ff45350572e7ba5fd0ff6185ab1251514e2f/src/main/java/com/gudos/persistence/Database.java\#L231}{\nolinkurl{persistence/Database.java}:231}; \href{https://github.com/itsnold/FiveNightsAtGudos/blob/e9f0ff45350572e7ba5fd0ff6185ab1251514e2f/src/main/java/com/gudos/persistence/SettingsStore.java\#L83}{\nolinkurl{persistence/SettingsStore.java}:83}; \href{https://github.com/itsnold/FiveNightsAtGudos/blob/e9f0ff45350572e7ba5fd0ff6185ab1251514e2f/src/main/java/com/gudos/persistence/ShiftLogStore.java\#L47}{\nolinkurl{persistence/ShiftLogStore.java}:47}}
\clearpage\pdfbookmark[0]{Patterns and implementation}{group-p:facade}\phantomsection\label{p:facade}\pdfbookmark[1]{Facade: New Game navigation and audio}{bm-p:facade}\section*{Facade: New Game navigation and audio}

\subsection*{Meaning and participants}
A facade exposes a small client operation while hiding coordination of a larger subsystem. Navigation is the screen-facing contract; Application is its concrete coordinator/facade. A method name or interface alone is not sufficient proof: the implementation must hide meaningful work.
\par\noindent\begin{tabularx}{\linewidth}{L{.25\linewidth}X}\toprule
\textbf{Client action}&\textbf{Request through the same boundary}\\\midrule
New Game&\code{navigation.startNight(1)}\\
Continue&\code{navigation.startNight(progress.nextNight())}\\
Select Night&\code{navigation.startNight(selected)}\\
Next / Play Again&\code{navigation.startNight(nextNight)}\\\bottomrule\end{tabularx}\par
\subsection*{What Application.startNight hides}
\begin{enumerate}
\item Guard busy/save/active-session/locked-night/incomplete-save cases.
\item Clear prior result state and capture the active profile name.
\item Construct GameSession with a new attempt identity.
\item Construct GamePanel with current-settings, sound and pause callbacks; attach hidden-game callback.
\item Replace the visible card, stop obsolete effects and request gameplay audio.
\end{enumerate}
The menu supplies intention, not window/storage/audio machinery. ``New Game'' is the label; a separate newGame method is unnecessary because starting Night 1 reuses validated night startup. Neither history deletion nor checkpoint restoration happens here.
\subsection*{A second facade over a different subsystem}
AudioService exposes play, ambience, menuMusic and gameplay. Its worker hides WAV/synthesis selection, stream decoding, Clip cache, volume/pan/mute controls, one-shots, loop replacement and cleanup. A client asks for a cue/layer, not a sequence of mixer calls.
\note{A concise explanation}{``The New Game button is a client of our navigation facade. Application.startNight hides the checks and construction needed to show a fresh playable night. AudioService separately hides Java Sound complexity behind cue requests.''}
\textbf{Tradeoff.} Central coordination prevents duplicated startup but gives Application many responsibilities. That is a practical coordinator design, not a claim that every method in it independently implements a pattern.
\sources{\href{https://github.com/itsnold/FiveNightsAtGudos/blob/e9f0ff45350572e7ba5fd0ff6185ab1251514e2f/src/main/java/com/gudos/ui/MenuPanel.java\#L86}{\nolinkurl{ui/MenuPanel.java}:86}; \href{https://github.com/itsnold/FiveNightsAtGudos/blob/e9f0ff45350572e7ba5fd0ff6185ab1251514e2f/src/main/java/com/gudos/ui/Navigation.java\#L46}{\nolinkurl{ui/Navigation.java}:46}; \href{https://github.com/itsnold/FiveNightsAtGudos/blob/e9f0ff45350572e7ba5fd0ff6185ab1251514e2f/src/main/java/com/gudos/Application.java\#L306}{\nolinkurl{Application.java}:306}; \href{https://github.com/itsnold/FiveNightsAtGudos/blob/e9f0ff45350572e7ba5fd0ff6185ab1251514e2f/src/main/java/com/gudos/audio/AudioService.java\#L122}{\nolinkurl{audio/AudioService.java}:122}}
\guidepage{Decorator: real composable camera effects}{p:decorator}

\par\noindent\begin{tabularx}{\linewidth}{L{.26\linewidth}X}\toprule\textbf{Pattern role}&\textbf{Project type}\\\midrule
Component&SceneRenderer interface\\Concrete component&CameraSceneRenderer\\Base decorator&RendererDecorator: is a SceneRenderer and has a wrapped SceneRenderer\\Concrete decorators&VhsRendererDecorator; DimmedRendererDecorator\\Composition client&GamePanel constructor\\\bottomrule\end{tabularx}\par
\begin{lstlisting}
SceneRenderer camera = new VhsRendererDecorator(
    new CameraSceneRenderer(session, assets), assets, .25f);
SceneRenderer switching = new VhsRendererDecorator(
    new DimmedRendererDecorator(camera, 115),
    assets, .85f, 3);
\end{lstlisting}
These locals illustrate the exact shapes used for GamePanel's two fields. Each wrapper delegates first, then adds its own layer. Normal draws room then light VHS. Switching draws room $\rightarrow$ normal VHS $\rightarrow$ darkness 115 $\rightarrow$ stronger outer VHS at triple speed.
\subsection*{Why this is Decorator}
The same render interface lets a wrapper replace a component and wrap another wrapper. It combines effects without one subclass for every combination. Layer order is observable: darkness after noise darkens noise; noise after darkness sits on top. The delegated drawing order determines those final pixels.
\subsection*{Details that support follow-up answers}
\begin{itemize}
\item RendererDecorator holds a nonnull SceneRenderer and delegates once with original time/reduced flag.
\item VhsRendererDecorator always draws the inner object, then skips optional noise for reduced mode/zero opacity. Its speed changes only its own effect time, not the inner scene or enemy rules.
\item DimmedRendererDecorator dims whenever called. GamePanel bypasses the switching chain in reduced mode; do not claim every wrapper individually checks the preference.
\item Effect drawing uses graphics.create() and finally/dispose to isolate wrapper-specific paint/composite/hints from later HUD drawing.
\end{itemize}
\note{What is not the same example}{VhsOverlay is a drawing utility, not itself a wrapper. BufferedInputStream is a Java-library decorator; the renderer classes are the project's custom implementation.}
\sources{\href{https://github.com/itsnold/FiveNightsAtGudos/blob/e9f0ff45350572e7ba5fd0ff6185ab1251514e2f/src/main/java/com/gudos/ui/rendering/RendererDecorator.java\#L40}{\nolinkurl{ui/rendering/RendererDecorator.java}:40}; \href{https://github.com/itsnold/FiveNightsAtGudos/blob/e9f0ff45350572e7ba5fd0ff6185ab1251514e2f/src/main/java/com/gudos/ui/rendering/VhsRendererDecorator.java\#L60}{\nolinkurl{ui/rendering/VhsRendererDecorator.java}:60}; \href{https://github.com/itsnold/FiveNightsAtGudos/blob/e9f0ff45350572e7ba5fd0ff6185ab1251514e2f/src/main/java/com/gudos/ui/rendering/DimmedRendererDecorator.java\#L43}{\nolinkurl{ui/rendering/DimmedRendererDecorator.java}:43}; \href{https://github.com/itsnold/FiveNightsAtGudos/blob/e9f0ff45350572e7ba5fd0ff6185ab1251514e2f/src/main/java/com/gudos/ui/GamePanel.java\#L98}{\nolinkurl{ui/GamePanel.java}:98}}
\guidepage{Builder, painting template and OOP choices}{p:builder}

\subsection*{Builder: create a level by named steps}
SecretLevelBuilder collects player/bundle spawn, destination, platforms, cover, patrols, width and countdown. Each setter returns this for chaining; repeated part calls append rather than replace. Application.createSecretLevel is the production recipe; build creates SecretLevel; SecretGame consumes that product.
\begin{lstlisting}
new SecretLevelBuilder().width(2560).player(80).bundle(240)
    .destination(new Rectangle(2410, 510, 125, 130))
    // repeated platform / cover / witness steps
    .timeLimit(60).build();
\end{lstlisting}
The recipe has three platforms, six cover regions and three independent patrols. build checks selected completeness/range conditions, not every possible NaN or geometry error. Construction copies input rectangles/lists, but returned Rectangle members remain mutable; the level is not deeply immutable after access.
\subsection*{Template Method style painting}
ScenePanel's final paintComponent fixes setup $\rightarrow$ transform/clip $\rightarrow$ subclass render $\rightarrow$ cleanup. paintChildren draws real controls then optional renderOverlay. Subclasses supply hooks without replacing the common transform skeleton. This is an additional behavioral structure; Builder is creational, while facade/decorator are structural.
\par\noindent\begin{tabularx}{\linewidth}{L{.24\linewidth}X}\toprule\textbf{OOP concept}&\textbf{Concrete evidence}\\\midrule
Encapsulation&Private GameSession state; guarded light/monitor methods.\\
Abstraction&Navigation and SceneRenderer small contracts.\\
Polymorphism&Current ScenePanel and plain/wrapped SceneRenderer share calls with different implementations.\\
Composition&GameSession owns GudoEncounter; decorator owns renderer.\\
Dependency injection&Main supplies dependencies; panels receive callbacks/settings supplier.\\
Value transfer&Settings, Progress and RunRecord records with generated accessors/equality.\\
Generics&submit<T> returns different worker values; place<T> preserves component type.\\\bottomrule\end{tabularx}\par
\note{Avoid overclaiming patterns}{One shared AssetManager is not enforced Singleton. Enum/conditional state machines are not separate polymorphic GoF State objects. Callbacks/listeners do not establish a separate custom Observer framework. Navigation is not adapting an incompatible interface.}
\sources{\href{https://github.com/itsnold/FiveNightsAtGudos/blob/e9f0ff45350572e7ba5fd0ff6185ab1251514e2f/src/main/java/com/gudos/Application.java\#L374}{\nolinkurl{Application.java}:374}; \href{https://github.com/itsnold/FiveNightsAtGudos/blob/e9f0ff45350572e7ba5fd0ff6185ab1251514e2f/src/main/java/com/gudos/game/minigame/SecretLevelBuilder.java\#L87}{\nolinkurl{game/minigame/SecretLevelBuilder.java}:87}; \href{https://github.com/itsnold/FiveNightsAtGudos/blob/e9f0ff45350572e7ba5fd0ff6185ab1251514e2f/src/main/java/com/gudos/ui/ScenePanel.java\#L116}{\nolinkurl{ui/ScenePanel.java}:116}}
\guidepage{Threads, audio layers and resource decisions}{p:services}

\par\noindent\begin{tabularx}{\linewidth}{L{.23\linewidth}X}\toprule\textbf{Context}&\textbf{Owned work}\\\midrule
Swing EDT&Window/screen changes, input callbacks, timer ticks and painting.\\
Boot SwingWorker&doInBackground loads files/assets; done installs UI on EDT.\\
gudo-storage&One serial worker for JDBC/settings/report operations; callbacks return to EDT.\\
gudo-audio&One serial daemon worker for decoding, mixer controls, playback and cleanup.\\\bottomrule\end{tabularx}\par
\subsection*{Why one game cue happens only once}
GamePanel compares location with lastGudoLocation for steps and office arrival cues. LEFT uses knock, RIGHT pound, FRONT whisper. At three attack seconds or less, rushPlayed gates urgency. A changed repelled count plays retreat. A one-shot warning occurs at power$\leq20$; room creaks/rattles have active-clock scheduling.
Watched arrival/departure disrupts only that viewed feed for two active seconds. Switching away silences it; returning before expiry preserves original variant/deadline. Mere selection of an already occupied feed does not create a sighting sting.
\subsection*{Playback implementation}
A cue is lazily decoded/cached as Clip. One-shots stop/reset/start the cached cue; repeated requests restart it rather than create unlimited voices. Layer requests compare sets, stopping removed loops and starting added ones while preserving shared positions. Master level times cue gain becomes dB through $20\log_{10}(v)$, clamped to supported controls. Approach pan maps by $2p-1$.
Missing mixer/cue failures enter unavailable and degrade to silence. Unsupported pan/gain controls are skipped. Only the three secret cues use Chiptune; current ordinary encounter cues use Sound-mapped WAVs. Archived source WAVs are not automatically active.
\subsection*{Resource and cleanup tradeoffs}
AssetManager caches decoded paths and requires clean defaults for all scenes. Unknown state names fall back; an explicitly listed missing asset fails boot. Buffers/blur are reused to reduce frame allocation, trading memory for work. Try-with-resources closes JDBC/streams; copied graphics contexts are disposed.
\note{Orderly exit}{Application stops timer/key dispatcher, queues audio cleanup, shuts down IO and disposes frame. Normal active-shift exit confirms and saves abandonment first; force-kill cannot execute this cleanup reliably.}
\sources{\href{https://github.com/itsnold/FiveNightsAtGudos/blob/e9f0ff45350572e7ba5fd0ff6185ab1251514e2f/src/main/java/com/gudos/Application.java\#L645}{\nolinkurl{Application.java}:645}; \href{https://github.com/itsnold/FiveNightsAtGudos/blob/e9f0ff45350572e7ba5fd0ff6185ab1251514e2f/src/main/java/com/gudos/ui/GamePanel.java\#L320}{\nolinkurl{ui/GamePanel.java}:320}; \href{https://github.com/itsnold/FiveNightsAtGudos/blob/e9f0ff45350572e7ba5fd0ff6185ab1251514e2f/src/main/java/com/gudos/audio/AudioService.java\#L183}{\nolinkurl{audio/AudioService.java}:183}; \href{https://github.com/itsnold/FiveNightsAtGudos/blob/e9f0ff45350572e7ba5fd0ff6185ab1251514e2f/src/main/java/com/gudos/assets/AssetManager.java\#L99}{\nolinkurl{assets/AssetManager.java}:99}}
\clearpage\pdfbookmark[0]{The three coordinating classes}{group-p:application}\phantomsection\label{p:application}\pdfbookmark[1]{Application.java: follow a request end to end}{bm-p:application}\section*{Application.java: follow a request end to end}

\phantomsection\label{file:Application}\pdfbookmark[2]{Application.java}{bm-file-Application}
\textbf{Role.} The Navigation implementation owns the current screen, live session/view, settings/progress, pending save and services. Panels ask for intent; this coordinator decides whether/how it can happen.
\par\noindent\begin{tabularx}{\linewidth}{L{.28\linewidth}X}\toprule\textbf{Method group}&\textbf{Order and reasoning}\\\midrule
constructor / dispatchKey&Receive boot dependencies; install close/focus/key handling; create timer; show menu. Ignore other-window events and debounce Escape/F11 before panel bindings.\\
tick / show&Measure/cap delta; update active normal model; handle ending; tick/repaint current screen; refresh audio. show replaces card and resets frame clock/focus.\\
menu / page / showMenu&Guard conflicting IO/active attempts. Resolve abandonment or incomplete saves; records needs background reads. Clearing a live UI reference does not delete saved progress.\\
startNight&Guard eligibility/pending state, capture player, create model and callback-driven view, attach secret, show and request audio. Reused by all normal start routes.\\
pauseShift / resumeShift&Release held controls and pause the same model. Resume reuses it; no new UUID or checkpoint is loaded.\\
openSecret / finishSecret&Pause normal session, create hidden model/view. Ending marks a secret outcome and retains the same normal night for restart after save.\\
finishShift / savePendingRun&Capture one final record/profile, choose result/scare, start IO. Keep UUID across failure; navigate/restart/quit only after the required save result.\\
saveSettings / display&Background preference/profile work, then replace UI values/apply display on EDT. F11 persists before changing native decorations; guard temporary focus loss.\\
quit / close / submit&Confirm/save abandonment on orderly exit. Stop timers/listeners/services. submit executes work on IO and posts success/failure to EDT.\\\bottomrule\end{tabularx}\par
\subsection*{State names worth recognizing}
\code{busy} blocks overlapping loads; \code{savingRun} blocks duplicate in-flight saves; \code{pendingRun/pendingProfile} survive retry; \code{activeProfile} fixes ownership at shift start. \code{quitAfterSave/menuAfterSave} defer navigation. \code{nextNightAfterSecret} is actually the same night. \code{displaySavePending/changingDisplay} guard native display changes.
\note{Trace this path aloud}{MenuPanel callback $\rightarrow$ startNight $\rightarrow$ GameSession/GamePanel $\rightarrow$ tick $\rightarrow$ finishShift $\rightarrow$ savePendingRun $\rightarrow$ Database then log $\rightarrow$ EDT result/navigation.}
\sources{\href{https://github.com/itsnold/FiveNightsAtGudos/blob/e9f0ff45350572e7ba5fd0ff6185ab1251514e2f/src/main/java/com/gudos/Application.java\#L306}{\nolinkurl{Application.java}:306}; \href{https://github.com/itsnold/FiveNightsAtGudos/blob/e9f0ff45350572e7ba5fd0ff6185ab1251514e2f/src/main/java/com/gudos/Application.java\#L454}{\nolinkurl{Application.java}:454}; \href{https://github.com/itsnold/FiveNightsAtGudos/blob/e9f0ff45350572e7ba5fd0ff6185ab1251514e2f/src/main/java/com/gudos/Application.java\#L486}{\nolinkurl{Application.java}:486}}
\guidepage{GameSession.java: rules and the final receipt}{p:gamesession}

\phantomsection\label{file:GameSession}\pdfbookmark[2]{GameSession.java}{bm-file-GameSession}
\textbf{Role.} One normal attempt with UUID/start Instant, night/duration, GudoEncounter, elapsed/power, monitor/light/pan/camera and optional scene overrides. A nonnull result is the finished marker.
\subsection*{Construction and update}
The one-argument constructor delegates with 360 s. The full constructor rejects night outside 1--5 and invalid duration. update validates delta, exits for pause/finish and processes exact deadlines as explained on page\pageref{p:rules}. No JFrame, audio service or database is needed to run these rules.
\par\noindent\begin{tabularx}{\linewidth}{L{.29\linewidth}X}\toprule\textbf{Method}&\textbf{Guard / effect / consequence}\\\midrule
\code{toggleMonitor}&Ignore paused/finished. Flip monitor and clear light. Lowering monitor does not automatically restore held light.\\
\code{setLight}&Always permit off; on requires active unfinished model and monitor closed. Correct facing is checked separately during update.\\
\code{selectCamera}&Validate camera1--6, then accept only with monitor up and active model. Closed-monitor shortcuts leave selection unchanged.\\
\code{setOfficePan}&Reject nonfinite input, clamp 0--1 during active play. Encounter uses this accepted direction.\\
\code{setSceneState / sceneState}&Store valid custom state. Actual singular enemy occupancy overrides it; otherwise use custom/default key. Images do not determine enemy logic.\\
\code{pause / resume}&Pause preserves memory and releases light; resume clears pause only if unfinished. Neither writes disk.\\
\code{finish / result}&Finish only once, freeze final values and equipment off. result throws before ending; later calls return the same receipt.\\
\code{abandon / secret endings}&Select distinct ABANDONED/MINIGAME outcome, not normal COMPLETED. Database progression depends on that distinction.\\
\code{clock / getters}&progress derives elapsed/duration. HUD clock is fictional; durationLabel is actual active time. Read accessors do not update/save.\\\bottomrule\end{tabularx}\par
\subsection*{Why immutable result identity matters}
Application can retry a failed storage sequence while presenting the same values and UUID. Creating a new receipt/UUID on retry would turn one attempt into multiple records. startedAt/endedAt are absolute timestamps, while elapsed excludes pause.
\note{What it intentionally cannot do}{There is no serialized current encounter/power/input checkpoint to restore after closing. Continue creates a new GameSession. The custom-duration constructor is a configuration facility, not a gameplay debug button.}
\sources{\href{https://github.com/itsnold/FiveNightsAtGudos/blob/e9f0ff45350572e7ba5fd0ff6185ab1251514e2f/src/main/java/com/gudos/game/GameSession.java\#L80}{\nolinkurl{game/GameSession.java}:80}; \href{https://github.com/itsnold/FiveNightsAtGudos/blob/e9f0ff45350572e7ba5fd0ff6185ab1251514e2f/src/main/java/com/gudos/game/GameSession.java\#L150}{\nolinkurl{game/GameSession.java}:150}; \href{https://github.com/itsnold/FiveNightsAtGudos/blob/e9f0ff45350572e7ba5fd0ff6185ab1251514e2f/src/main/java/com/gudos/game/GameSession.java\#L271}{\nolinkurl{game/GameSession.java}:271}}
\guidepage{GamePanel.java: input, view and accepted feedback}{p:gamepanel}

\phantomsection\label{file:GamePanel}\pdfbookmark[2]{GamePanel.java}{bm-file-GamePanel}
\textbf{Role.} This ScenePanel receives a model, shared assets, a current-settings supplier, sound callback and pause callback. It owns screen/input animation state, not authoritative save/progression rules.
\subsection*{Follow a button or key}
Constructor builds real monitor/light/camera GameButtons and press/release bindings. Space/Tab have held flags to avoid repeat toggles. F and mouse light combine into one requested light state; releaseControls clears all held flags on screen changes. The model still rejects illegal device combinations.
\par\noindent\begin{tabularx}{\linewidth}{L{.29\linewidth}X}\toprule\textbf{Method group}&\textbf{Responsibility}\\\midrule
\code{updateLight / toggleMonitor / selectCamera}&Ask the model first; emit audio/transition only for a real accepted change. Reselecting current camera or blocked F stays silent.\\
\code{syncControls}&Mirror model state into button visibility, selection, label and accessibility. Camera buttons appear with monitor up; light appears in office.\\
\code{tick}&After model update, ease office input pan and send it back. Detect new locations/repelled for one-time cues, feed disruption, low-power/urgency and scheduled room noises.\\
\code{render / renderHud}&Choose office/CCTV then draw flat clock, night, power, usage, approach instructions and fade progress. Drawing does not advance the model.\\
\code{renderOffice}&Blend separate clean/occupied projection outputs during exposure, then light/tint/vignette and readability shading.\\
\code{renderCamera}&Obscure a disrupted feed, otherwise choose normal/switching renderer chain. Add border, record dot, room label and schematic map.\\
\code{mouseClicked / onSecret}&Reverse-map five valid CAM04 paper clicks and run Application's supplied callback. The panel never constructs the hidden level itself.\\
\code{renderOverlay}&After Swing children, add optional subtle static and intro; reduced mode bypasses motion.\\\bottomrule\end{tabularx}\par
\subsection*{Presentation state to recognize}
\code{targetPan/pan} are eased office direction; separate projection buffers support blending. \code{lastGudoLocation/lastRepelled} gate cues. \code{garbledCamera/garbleUntil/garbleSound} retain one feed's disruption deadline/variant. \code{rushPlayed/lowPowerPlayed} prevent repeated warnings. \code{paperClicks} counts valid hits.
\note{A subtle update detail}{Application updates GameSession before GamePanel.tick sends its newest pan back, so direction can have one small update-step delay. All of this occurs on EDT; it is not a separate model/render thread.}
\sources{\href{https://github.com/itsnold/FiveNightsAtGudos/blob/e9f0ff45350572e7ba5fd0ff6185ab1251514e2f/src/main/java/com/gudos/ui/GamePanel.java\#L320}{\nolinkurl{ui/GamePanel.java}:320}; \href{https://github.com/itsnold/FiveNightsAtGudos/blob/e9f0ff45350572e7ba5fd0ff6185ab1251514e2f/src/main/java/com/gudos/ui/GamePanel.java\#L423}{\nolinkurl{ui/GamePanel.java}:423}; \href{https://github.com/itsnold/FiveNightsAtGudos/blob/e9f0ff45350572e7ba5fd0ff6185ab1251514e2f/src/main/java/com/gudos/ui/GamePanel.java\#L183}{\nolinkurl{ui/GamePanel.java}:183}}
\clearpage\pdfbookmark[0]{Every production Java file}{group-p:files-1}\phantomsection\label{p:files-1}\pdfbookmark[1]{Startup and screen requests}{bm-p:files-1}\section*{Startup and screen requests}
{\small Each card gives the purpose, execution path and reason. Detailed per-method explanations remain immediately above the actual source methods.}\par\vspace{2pt}\filecard{Main}{\href{https://github.com/itsnold/FiveNightsAtGudos/blob/e9f0ff45350572e7ba5fd0ff6185ab1251514e2f/src/main/java/com/gudos/Main.java\#L31}{\nolinkurl{Main.java}:31}}{Starts the desktop application and loads its dependencies.\par{\sffamily\scriptsize\color{Muted}\textbf{Find:} main, launch, SwingWorker.doInBackground/done, icon}}{main schedules launch on Swing; launch shows LoadingPanel; SwingWorker loads paths, settings, database, progress and images; done constructs Application.}{The window can repaint while disk and image loading happen. Boot is a record that transfers the loaded values together. Startup failures display a message and dispose the window.}
\filecard{Navigation}{\href{https://github.com/itsnold/FiveNightsAtGudos/blob/e9f0ff45350572e7ba5fd0ff6185ab1251514e2f/src/main/java/com/gudos/ui/Navigation.java\#L24}{\nolinkurl{ui/Navigation.java}:24}}{Defines the small screen-facing facade contract implemented by Application.\par{\sffamily\scriptsize\color{Muted}\textbf{Find:} menu, page, startNight, saveSettings, resumeShift, leaveShift, quit}}{MenuPanel, ExtrasPanel, NightSelectPanel, SettingsPanel, PausePanel and ResultPanel depend on this interface.}{Screens request intent without knowing JFrame, storage workers, GameSession construction or AudioService internals. An interface alone is not proof of Facade; the hidden orchestration in Application supplies the evidence.}
\filecard{MenuPanel}{\href{https://github.com/itsnold/FiveNightsAtGudos/blob/e9f0ff45350572e7ba5fd0ff6185ab1251514e2f/src/main/java/com/gudos/ui/MenuPanel.java\#L33}{\nolinkurl{ui/MenuPanel.java}:33}}{Displays the title menu, enabled Continue state and random portrait flicker.\par{\sffamily\scriptsize\color{Muted}\textbf{Find:} constructor, entry, move/select, tick, render/renderOverlay}}{New Game -> startNight(1); Continue -> startNight(progress.nextNight()); Extras -> page(EXTRAS); Quit -> quit.}{New Game deliberately preserves existing history. Continue is enabled after any recorded attempt, even a loss. Keyboard selection skips disabled entries. Flicker randomness is visual only.}

\guidepage{Extras and static screens}{p:files-2}
{\small Each card gives the purpose, execution path and reason. Detailed per-method explanations remain immediately above the actual source methods.}\par\vspace{2pt}\filecard{ExtrasPanel}{\href{https://github.com/itsnold/FiveNightsAtGudos/blob/e9f0ff45350572e7ba5fd0ff6185ab1251514e2f/src/main/java/com/gudos/ui/ExtrasPanel.java\#L23}{\nolinkurl{ui/ExtrasPanel.java}:23}}{Routes five text entries to secondary pages.\par{\sffamily\scriptsize\color{Muted}\textbf{Find:} constructor page callbacks; render}}{Select Night, How to Play, Shift Records, Settings and Credits map in order to Navigation.Page values.}{Keep a small routing screen without placing database or gameplay logic here. Capturing index per loop keeps each callback attached to the correct page.}
\filecard{NightSelectPanel}{\href{https://github.com/itsnold/FiveNightsAtGudos/blob/e9f0ff45350572e7ba5fd0ff6185ab1251514e2f/src/main/java/com/gudos/ui/NightSelectPanel.java\#L25}{\nolinkurl{ui/NightSelectPanel.java}:25}}{Shows five night buttons gated by the loaded progress.\par{\sffamily\scriptsize\color{Muted}\textbf{Find:} constructor eligibility/callbacks; render}}{Extras -> Select Night -> enabled Night N -> navigation.startNight(N).}{UI disables locked nights; Application rechecks unlocked before starting. Capture the loop value in selected so each listener keeps its own night.}
\filecard{InfoPanel}{\href{https://github.com/itsnold/FiveNightsAtGudos/blob/e9f0ff45350572e7ba5fd0ff6185ab1251514e2f/src/main/java/com/gudos/ui/InfoPanel.java\#L25}{\nolinkurl{ui/InfoPanel.java}:25}}{Draws the in-game guide or credits.\par{\sffamily\scriptsize\color{Muted}\textbf{Find:} render, renderGuide, renderCredits}}{Application.page(GUIDE/CREDITS) constructs the same class with credits=false/true.}{Static information can share one view and backdrop. Guide controls are descriptive text; GamePanel/ScenePanel implement the actual actions.}

\guidepage{Shared Swing infrastructure}{p:files-3}
{\small Each card gives the purpose, execution path and reason. Detailed per-method explanations remain immediately above the actual source methods.}\par\vspace{2pt}\filecard{ScenePanel}{\href{https://github.com/itsnold/FiveNightsAtGudos/blob/e9f0ff45350572e7ba5fd0ff6185ab1251514e2f/src/main/java/com/gudos/ui/ScenePanel.java\#L30}{\nolinkurl{ui/ScenePanel.java}:30}}{Provides a common scaled scene, real Swing component layout and keyboard routing.\par{\sffamily\scriptsize\color{Muted}\textbf{Find:} place, doLayout, scenePoint, paintComponent/paintChildren, bind/dispatchKey}}{Subclasses place controls in 1280x720 coordinates and implement render; Swing invokes paintComponent/paintChildren/doLayout.}{Uniform scale and centered offsets preserve aspect ratio. Reverse transforms preserve input accuracy. The final paint skeleton is a Template Method style design with subclass hooks.}
\filecard{BackdropPanel}{\href{https://github.com/itsnold/FiveNightsAtGudos/blob/e9f0ff45350572e7ba5fd0ff6185ab1251514e2f/src/main/java/com/gudos/ui/BackdropPanel.java\#L26}{\nolinkurl{ui/BackdropPanel.java}:26}}{Shares the dim office background and navigation chrome for secondary pages.\par{\sffamily\scriptsize\color{Muted}\textbf{Find:} backdrop, pageChrome, backButton}}{Extras, settings, guide/credits, selection, pause, records and result reuse backdrop or pageChrome.}{Avoid duplicating the same drawing and Back binding. It is an abstract base class; inheritance alone is not the custom Decorator pattern.}
\filecard{GameButton}{\href{https://github.com/itsnold/FiveNightsAtGudos/blob/e9f0ff45350572e7ba5fd0ff6185ab1251514e2f/src/main/java/com/gudos/ui/GameButton.java\#L24}{\nolinkurl{ui/GameButton.java}:24}}{Uses a real JButton with custom visual styles and a supplied action.\par{\sffamily\scriptsize\color{Muted}\textbf{Find:} constructor listener; paintComponent}}{Screens construct the button with label/style/callback; Swing fires ActionListener; the callback requests navigation or model input.}{Preserve enabled state, keyboard/button behavior and accessibility while drawing the game style. Painting only changes appearance, not game rules.}

\guidepage{Common drawing and loading}{p:files-4}
{\small Each card gives the purpose, execution path and reason. Detailed per-method explanations remain immediately above the actual source methods.}\par\vspace{2pt}\filecard{Theme}{\href{https://github.com/itsnold/FiveNightsAtGudos/blob/e9f0ff45350572e7ba5fd0ff6185ab1251514e2f/src/main/java/com/gudos/ui/Theme.java\#L25}{\nolinkurl{ui/Theme.java}:25}}{Centralizes the bundled pixel font, colors and reusable drawing utilities.\par{\sffamily\scriptsize\color{Muted}\textbf{Find:} font helpers, text/paragraph/ellipsize, cover/contain, quality/vignette}}{Panels/buttons/renderers call static helpers instead of reproducing font/image/text code.}{One style source keeps UI consistent. Explicit wrap/ellipsize handle Java2D text limits. cover crops to fill while contain preserves the entire image.}
\filecard{VhsOverlay}{\href{https://github.com/itsnold/FiveNightsAtGudos/blob/e9f0ff45350572e7ba5fd0ff6185ab1251514e2f/src/main/java/com/gudos/ui/VhsOverlay.java\#L25}{\nolinkurl{ui/VhsOverlay.java}:25}}{Draws animated noise, scanlines and a rolling band as a visual utility.\par{\sffamily\scriptsize\color{Muted}\textbf{Find:} paint}}{Menu/GamePanel and VhsRendererDecorator call paint.}{It returns immediately for reduced effects and restores composite/interpolation. The decorator uses a copied Graphics2D so color/paint changes do not leak into later HUD drawing.}
\filecard{LoadingPanel}{\href{https://github.com/itsnold/FiveNightsAtGudos/blob/e9f0ff45350572e7ba5fd0ff6185ab1251514e2f/src/main/java/com/gudos/ui/LoadingPanel.java\#L22}{\nolinkurl{ui/LoadingPanel.java}:22}}{Draws a waiting message and animated dots.\par{\sffamily\scriptsize\color{Muted}\textbf{Find:} constructor; render}}{Main uses it during boot; Application uses it for asynchronous records/progress reads.}{The panel does not perform slow work. The worker and coordinator decide when loading is complete.}

\guidepage{Pause, result and scare screens}{p:files-5}
{\small Each card gives the purpose, execution path and reason. Detailed per-method explanations remain immediately above the actual source methods.}\par\vspace{2pt}\filecard{PausePanel}{\href{https://github.com/itsnold/FiveNightsAtGudos/blob/e9f0ff45350572e7ba5fd0ff6185ab1251514e2f/src/main/java/com/gudos/ui/PausePanel.java\#L25}{\nolinkurl{ui/PausePanel.java}:25}}{Displays a frozen active attempt and routes Resume or Return to Menu.\par{\sffamily\scriptsize\color{Muted}\textbf{Find:} constructor callbacks; render}}{Escape in gameplay -> Application.pauseShift -> PausePanel; Resume/Escape -> resumeShift; Return to Menu -> leaveShift confirmation/save.}{Reuse the same model when resuming. Leaving creates ABANDONED, so pause is different from completing an attempt.}
\filecard{ResultPanel}{\href{https://github.com/itsnold/FiveNightsAtGudos/blob/e9f0ff45350572e7ba5fd0ff6185ab1251514e2f/src/main/java/com/gudos/ui/ResultPanel.java\#L28}{\nolinkurl{ui/ResultPanel.java}:28}}{Displays an ended RunRecord and the pending/saved/failed storage status.\par{\sffamily\scriptsize\color{Muted}\textbf{Find:} saving, saved, failed, render}}{Application.finishShift builds it; savePendingRun drives saving/saved/failed; buttons request next/replay/menu/retry.}{Disable next/menu while saving; allow retry with the same pending UUID. A completed Night 1-4 offers next; other endings offer the same night; Night 5 remains capped.}
\filecard{JumpscarePanel}{\href{https://github.com/itsnold/FiveNightsAtGudos/blob/e9f0ff45350572e7ba5fd0ff6185ab1251514e2f/src/main/java/com/gudos/ui/JumpscarePanel.java\#L29}{\nolinkurl{ui/JumpscarePanel.java}:29}}{Animates the caught still and calls back once after 2.1 seconds.\par{\sffamily\scriptsize\color{Muted}\textbf{Find:} constructor, tick, render}}{Application.finishShift on CAUGHT starts sound/save and shows this; completion opens the prepared ResultPanel.}{Precompute a horizontal blur once. Zoom/shake reuse the still; reducedEffects keeps it stationary. Guard delivered avoids repeated navigation.}

\guidepage{Records and preference views}{p:files-6}
{\small Each card gives the purpose, execution path and reason. Detailed per-method explanations remain immediately above the actual source methods.}\par\vspace{2pt}\filecard{RecordsPanel}{\href{https://github.com/itsnold/FiveNightsAtGudos/blob/e9f0ff45350572e7ba5fd0ff6185ab1251514e2f/src/main/java/com/gudos/ui/RecordsPanel.java\#L35}{\nolinkurl{ui/RecordsPanel.java}:35}}{Displays the current profile summary and read-only recent attempts table.\par{\sffamily\scriptsize\color{Muted}\textbf{Find:} constructor/table model; render}}{Extras -> Shift Records -> Application IO load -> constructor builds JTable from the returned list.}{No SQL is run while rendering. Absolute Instants are displayed in the local system zone. Only latest 200 rows are loaded; summary counts include all stored attempts.}
\filecard{SettingsPanel}{\href{https://github.com/itsnold/FiveNightsAtGudos/blob/e9f0ff45350572e7ba5fd0ff6185ab1251514e2f/src/main/java/com/gudos/ui/SettingsPanel.java\#L39}{\nolinkurl{ui/SettingsPanel.java}:39}}{Edits preferences in real Swing controls and submits an immutable Settings value.\par{\sffamily\scriptsize\color{Muted}\textbf{Find:} values, setSaving/setFullscreen, DocumentFilter, checkbox icons}}{Extras -> Settings -> edit name/volume/effects/fullscreen -> Apply -> navigation.saveSettings(values()).}{DocumentFilter sanitizes typed/pasted names, and Settings validates again. Applying is explicit; Back does not save edits. The clipboard button copies the resolved data folder.}
\filecard{Settings}{\href{https://github.com/itsnold/FiveNightsAtGudos/blob/e9f0ff45350572e7ba5fd0ff6185ab1251514e2f/src/main/java/com/gudos/persistence/Settings.java\#L21}{\nolinkurl{persistence/Settings.java}:21}}{Validates and carries immutable user preferences.\par{\sffamily\scriptsize\color{Muted}\textbf{Find:} compact constructor; withFullscreen}}{SettingsStore loads it; SettingsPanel creates a replacement; Application persists and applies it.}{Sanitize control characters, trim/fallback/cap the name, and clamp volume at the value boundary as well as the UI. A record replacement makes updates explicit.}

\guidepage{Data paths and value carriers}{p:files-7}
{\small Each card gives the purpose, execution path and reason. Detailed per-method explanations remain immediately above the actual source methods.}\par\vspace{2pt}\filecard{AppPaths}{\href{https://github.com/itsnold/FiveNightsAtGudos/blob/e9f0ff45350572e7ba5fd0ff6185ab1251514e2f/src/main/java/com/gudos/persistence/AppPaths.java\#L26}{\nolinkurl{persistence/AppPaths.java}:26}}{Resolves writable player-data paths outside the packaged game.\par{\sffamily\scriptsize\color{Muted}\textbf{Find:} resolve; database/settings/shiftLog}}{Main calls resolve; SettingsPanel displays the folder; Database, SettingsStore and ShiftLogStore use the three child paths.}{A JAR contains read-only resources; user saves need a writable external folder. gudo.dataDir supports isolated demo runs without changing ordinary player data.}
\filecard{Progress}{\href{https://github.com/itsnold/FiveNightsAtGudos/blob/e9f0ff45350572e7ba5fd0ff6185ab1251514e2f/src/main/java/com/gudos/persistence/Progress.java\#L21}{\nolinkurl{persistence/Progress.java}:21}}{Provides a loaded summary for menu, night-selection and records screens.\par{\sffamily\scriptsize\color{Muted}\textbf{Find:} EMPTY; unlocked}}{Database.loadProgress calculates counts and reads unlock fields; Application passes the value to panels.}{Keep one summarized value for UI decisions instead of running SQL during every repaint. EMPTY represents a fresh profile.}
\filecard{RunRecord}{\href{https://github.com/itsnold/FiveNightsAtGudos/blob/e9f0ff45350572e7ba5fd0ff6185ab1251514e2f/src/main/java/com/gudos/game/RunRecord.java\#L23}{\nolinkurl{game/RunRecord.java}:23}}{Carries the immutable snapshot of one ended attempt.\par{\sffamily\scriptsize\color{Muted}\textbf{Find:} generated accessors; durationLabel}}{GameSession creates it once; Application retains it; Database and ShiftLogStore save the same ID and values; RecordsPanel displays it.}{A record supplies value equality and accessors. Capture final values once so retries cannot silently turn the attempt into a new one.}

\guidepage{Named states and camera motion}{p:files-8}
{\small Each card gives the purpose, execution path and reason. Detailed per-method explanations remain immediately above the actual source methods.}\par\vspace{2pt}\filecard{Outcome}{\href{https://github.com/itsnold/FiveNightsAtGudos/blob/e9f0ff45350572e7ba5fd0ff6185ab1251514e2f/src/main/java/com/gudos/game/Outcome.java\#L20}{\nolinkurl{game/Outcome.java}:20}}{Defines the six final result categories stored in SQLite.\par{\sffamily\scriptsize\color{Muted}\textbf{Find:} stable enum name; label}}{GameSession.finish selects the enum; Database stores name(); UI/logs display label().}{Stable machine names separate storage from readable labels. Only COMPLETED advances normal-night progress; MINIGAME\_COMPLETED is a different outcome.}
\filecard{Scene}{\href{https://github.com/itsnold/FiveNightsAtGudos/blob/e9f0ff45350572e7ba5fd0ff6185ab1251514e2f/src/main/java/com/gudos/game/Scene.java\#L22}{\nolinkurl{game/Scene.java}:22}}{Defines the office and six canonical camera rooms.\par{\sffamily\scriptsize\color{Muted}\textbf{Find:} id/label/camera; static camera(number)}}{Game controls validate a camera number with Scene.camera(number); logic and AssetManager share the enum id.}{Central identifiers keep camera labels, route locations and image manifest keys aligned. CAM 01 is currently Front entrance, despite some older dining filenames/text.}
\filecard{CameraMotion}{\href{https://github.com/itsnold/FiveNightsAtGudos/blob/e9f0ff45350572e7ba5fd0ff6185ab1251514e2f/src/main/java/com/gudos/game/CameraMotion.java\#L21}{\nolinkurl{game/CameraMotion.java}:21}}{Computes the automatic CCTV sweep from camera identity and active shift time.\par{\sffamily\scriptsize\color{Muted}\textbf{Find:} pan}}{CameraSceneRenderer and the secret-paper hit test call pan(scene,elapsed).}{A cosine gives smooth reversal. Each camera has a different phase and period. Using elapsed means unseen cameras continue their sweep and all sweeps freeze during pause.}

\guidepage{Image loading and projection}{p:files-9}
{\small Each card gives the purpose, execution path and reason. Detailed per-method explanations remain immediately above the actual source methods.}\par\vspace{2pt}\filecard{AssetManager}{\href{https://github.com/itsnold/FiveNightsAtGudos/blob/e9f0ff45350572e7ba5fd0ff6185ab1251514e2f/src/main/java/com/gudos/assets/AssetManager.java\#L35}{\nolinkurl{assets/AssetManager.java}:35}}{Loads bundled images, the scene manifest and decoded TV-noise frames, then serves cached lookups.\par{\sffamily\scriptsize\color{Muted}\textbf{Find:} constructor, resource, scene, noise/menuRender}}{Main constructs one manager in the boot worker. Screens ask scene(scene,state), menuRender(index) or noise(frame); resource opens a classpath stream.}{Decode each referenced path once, require a default for all rooms, and fail startup for a listed missing asset. Unknown state names fall back to a clean default. This separates a bad key from a broken packaged asset.}
\filecard{PanoramaRenderer}{\href{https://github.com/itsnold/FiveNightsAtGudos/blob/e9f0ff45350572e7ba5fd0ff6185ab1251514e2f/src/main/java/com/gudos/assets/PanoramaRenderer.java\#L33}{\nolinkurl{assets/PanoramaRenderer.java}:33}}{Projects a wide room image into a curved-looking view and maps clicks back to source pixels.\par{\sffamily\scriptsize\color{Muted}\textbf{Find:} prepare, render, sourcePoint}}{render clamps/quantizes pan, prepares sampling maps when the source changes, fills a reused 854x480 image, and returns it. sourcePoint reverses the latest sampling for a hit test.}{atan maps horizontal rays and cos changes vertical sampling near the edges. Cached pixel maps and one output buffer reduce allocation. A separate renderer is needed when two outputs must coexist for a blend.}
\filecard{CameraSceneRenderer}{\href{https://github.com/itsnold/FiveNightsAtGudos/blob/e9f0ff45350572e7ba5fd0ff6185ab1251514e2f/src/main/java/com/gudos/ui/rendering/CameraSceneRenderer.java\#L32}{\nolinkurl{ui/rendering/CameraSceneRenderer.java}:32}}{Draws the current CCTV room without VHS or switching effects.\par{\sffamily\scriptsize\color{Muted}\textbf{Find:} constructor; render}}{GamePanel puts it inside VhsRendererDecorator; render reads selected camera, room state and sweep from the live session.}{Composition keeps one core room renderer separate from optional visual layers. The object is created once and reads updated model values each frame.}

\guidepage{Renderer contracts and wrappers}{p:files-10}
{\small Each card gives the purpose, execution path and reason. Detailed per-method explanations remain immediately above the actual source methods.}\par\vspace{2pt}\filecard{SceneRenderer}{\href{https://github.com/itsnold/FiveNightsAtGudos/blob/e9f0ff45350572e7ba5fd0ff6185ab1251514e2f/src/main/java/com/gudos/ui/rendering/SceneRenderer.java\#L24}{\nolinkurl{ui/rendering/SceneRenderer.java}:24}}{Defines the component interface shared by a room renderer and all custom renderer decorators.\par{\sffamily\scriptsize\color{Muted}\textbf{Find:} render contract}}{GamePanel stores composed implementations through SceneRenderer and calls render(graphics,time,reducedEffects).}{The same interface allows nested wrappers and substitutability. Without it, effect layers would need to know each concrete renderer type.}
\filecard{RendererDecorator}{\href{https://github.com/itsnold/FiveNightsAtGudos/blob/e9f0ff45350572e7ba5fd0ff6185ab1251514e2f/src/main/java/com/gudos/ui/rendering/RendererDecorator.java\#L24}{\nolinkurl{ui/rendering/RendererDecorator.java}:24}}{Implements SceneRenderer while retaining another SceneRenderer and delegating to it.\par{\sffamily\scriptsize\color{Muted}\textbf{Find:} constructor; render delegation}}{Concrete effect constructors call super(wrapped); render forwards the original graphics/time/reducedEffects.}{Being the same component type and containing one component lets decorators wrap decorators. Null rejection fails invalid chains early.}
\filecard{VhsRendererDecorator}{\href{https://github.com/itsnold/FiveNightsAtGudos/blob/e9f0ff45350572e7ba5fd0ff6185ab1251514e2f/src/main/java/com/gudos/ui/rendering/VhsRendererDecorator.java\#L27}{\nolinkurl{ui/rendering/VhsRendererDecorator.java}:27}}{Adds optional VHS interference after its wrapped renderer draws.\par{\sffamily\scriptsize\color{Muted}\textbf{Find:} constructors; render}}{Normal camera uses opacity 0.25; switching adds an outer opacity 0.85 layer at speed 3.}{Delegate with original time so effect speed affects only this wrapper. Skip noise when reducedEffects or opacity=0, but still render the wrapped room.}

\guidepage{Darkness and audio services}{p:files-11}
{\small Each card gives the purpose, execution path and reason. Detailed per-method explanations remain immediately above the actual source methods.}\par\vspace{2pt}\filecard{DimmedRendererDecorator}{\href{https://github.com/itsnold/FiveNightsAtGudos/blob/e9f0ff45350572e7ba5fd0ff6185ab1251514e2f/src/main/java/com/gudos/ui/rendering/DimmedRendererDecorator.java\#L25}{\nolinkurl{ui/rendering/DimmedRendererDecorator.java}:25}}{Adds a translucent black layer after its wrapped renderer draws.\par{\sffamily\scriptsize\color{Muted}\textbf{Find:} constructor; render}}{GamePanel puts it around the normal camera chain for switching interference.}{Validate alpha=0-255, delegate first, and draw using a copied graphics context. Layer order changes pixels.}
\filecard{AudioService}{\href{https://github.com/itsnold/FiveNightsAtGudos/blob/e9f0ff45350572e7ba5fd0ff6185ab1251514e2f/src/main/java/com/gudos/audio/AudioService.java\#L39}{\nolinkurl{audio/AudioService.java}:39}}{Provides the audio facade over Java Sound and a serial audio worker.\par{\sffamily\scriptsize\color{Muted}\textbf{Find:} clip, play, ambience/gameplay, gain/pan, stop/close}}{Application or a panel callback requests play, ambience, menuMusic or gameplay. The worker lazily decodes/caches a Clip, sets gain/pan, starts or loops it, and later stops/closes it.}{One facade hides stream decoding, mixer access, caching, thread scheduling, layers and cleanup. Missing hardware is remembered as unavailable and gameplay continues. The worker is daemon; the storage worker is not.}
\filecard{Sound}{\href{https://github.com/itsnold/FiveNightsAtGudos/blob/e9f0ff45350572e7ba5fd0ff6185ab1251514e2f/src/main/java/com/gudos/audio/Sound.java\#L21}{\nolinkurl{audio/Sound.java}:21}}{Names each cue and stores its classpath path and relative gain.\par{\sffamily\scriptsize\color{Muted}\textbf{Find:} enum path/gain accessors}}{GamePanel/SecretPanel/Application choose a Sound enum; AudioService reads path and gain.}{Enum values avoid scattered filename strings. Master volume and per-sound gain are separate. Null paths mark the three Chiptune sounds.}

\guidepage{Synthesis and text stores}{p:files-12}
{\small Each card gives the purpose, execution path and reason. Detailed per-method explanations remain immediately above the actual source methods.}\par\vspace{2pt}\filecard{Chiptune}{\href{https://github.com/itsnold/FiveNightsAtGudos/blob/e9f0ff45350572e7ba5fd0ff6185ab1251514e2f/src/main/java/com/gudos/audio/Chiptune.java\#L26}{\nolinkurl{audio/Chiptune.java}:26}}{Generates only the hidden minigame audio in memory.\par{\sffamily\scriptsize\color{Muted}\textbf{Find:} create; ambience}}{create returns an AudioInputStream for SECRET\_JUMP, SECRET\_PICKUP or SECRET\_AMBIENCE. AudioService opens it through the same clip path used for WAVs.}{Square waves and envelopes create small retro cues without another library. The eight-second drone uses a deterministic noise sequence and quiet endpoints. Current ordinary encounter sounds are bundled WAVs.}
\filecard{SettingsStore}{\href{https://github.com/itsnold/FiveNightsAtGudos/blob/e9f0ff45350572e7ba5fd0ff6185ab1251514e2f/src/main/java/com/gudos/persistence/SettingsStore.java\#L32}{\nolinkurl{persistence/SettingsStore.java}:32}}{Reads and replaces the UTF-8 settings.txt preferences file.\par{\sffamily\scriptsize\color{Muted}\textbf{Find:} load, booleanValue, save}}{load creates defaults on first launch or parses Properties. save writes a temporary neighbor, then replaces the real file.}{Use explicit true/false parsing, volume fallback and Settings normalization. Atomic replacement reduces truncated settings; filesystems without atomic move use a regular replacement fallback.}
\filecard{ShiftLogStore}{\href{https://github.com/itsnold/FiveNightsAtGudos/blob/e9f0ff45350572e7ba5fd0ff6185ab1251514e2f/src/main/java/com/gudos/persistence/ShiftLogStore.java\#L32}{\nolinkurl{persistence/ShiftLogStore.java}:32}}{Appends a readable one-line report for each ended attempt.\par{\sffamily\scriptsize\color{Muted}\textbf{Find:} append}}{Application.savePendingRun calls append after Database.saveRun.}{A run=UUID marker detects an already-written report on retry. Locale.ROOT stabilizes decimal formatting; UTF-8 preserves names. This is best-effort application-level deduplication, not a transactional database.}

\guidepage{Database and level construction}{p:files-13}
{\small Each card gives the purpose, execution path and reason. Detailed per-method explanations remain immediately above the actual source methods.}\par\vspace{2pt}\filecard{Database}{\href{https://github.com/itsnold/FiveNightsAtGudos/blob/e9f0ff45350572e7ba5fd0ff6185ab1251514e2f/src/main/java/com/gudos/persistence/Database.java\#L38}{\nolinkurl{persistence/Database.java}:38}}{Persists local profiles, unlock summaries and attempt history using SQLite/JDBC.\par{\sffamily\scriptsize\color{Muted}\textbf{Find:} constructor/migrations, connect/profile, loadProgress, saveRun/loadRuns}}{Constructor configures/migrates schema; loadProgress selects/creates a profile; saveRun inserts an attempt and updates unlocks transactionally; loadRuns returns latest rows.}{Prepared statements bind player values. Foreign keys and CHECK/UNIQUE constraints protect relations and ranges. A UUID primary key supports idempotent retries; MAX preserves unlocks when replaying earlier nights.}
\filecard{SecretLevelBuilder}{\href{https://github.com/itsnold/FiveNightsAtGudos/blob/e9f0ff45350572e7ba5fd0ff6185ab1251514e2f/src/main/java/com/gudos/game/minigame/SecretLevelBuilder.java\#L30}{\nolinkurl{game/minigame/SecretLevelBuilder.java}:30}}{Constructs a SecretLevel step by step through a fluent Builder.\par{\sffamily\scriptsize\color{Muted}\textbf{Find:} fluent part methods; build}}{Application.createSecretLevel chains width/player/bundle/destination/platform/cover/witness/timeLimit, then build.}{Optional repeated parts are easier to read than a large positional constructor. Each setter returns this; build validates completeness and creates the product. Validation covers selected constraints, not every possible numeric input.}
\filecard{SecretLevel}{\href{https://github.com/itsnold/FiveNightsAtGudos/blob/e9f0ff45350572e7ba5fd0ff6185ab1251514e2f/src/main/java/com/gudos/game/minigame/SecretLevel.java\#L25}{\nolinkurl{game/minigame/SecretLevel.java}:25}}{Stores the level configuration produced by SecretLevelBuilder.\par{\sffamily\scriptsize\color{Muted}\textbf{Find:} compact constructor; Patrol}}{Application creates the product; SecretGame reads dimensions, spawn points, cover, platforms, destination and patrols.}{The constructor copies incoming Rectangle objects and lists, so later changes to builder inputs do not change this product. Returned Rectangles remain mutable; this is not deep immutability after access.}

\guidepage{Enemy and hidden-game runtime}{p:files-14}
{\small Each card gives the purpose, execution path and reason. Detailed per-method explanations remain immediately above the actual source methods.}\par\vspace{2pt}\filecard{GudoEncounter}{\href{https://github.com/itsnold/FiveNightsAtGudos/blob/e9f0ff45350572e7ba5fd0ff6185ab1251514e2f/src/main/java/com/gudos/game/GudoEncounter.java\#L23}{\nolinkurl{game/GudoEncounter.java}:23}}{Implements one deterministic enemy encounter state machine and three approach routes.\par{\sffamily\scriptsize\color{Muted}\textbf{Find:} untilEvent, advance, location, faces, exposure/timer accessors}}{GameSession asks untilEvent, advances that exact step, then checks caught. The UI reads location, approach, exposure/fade and attack countdown.}{routeIndex=-1 represents waiting, camera indices represent travel, atOffice represents defense, and caught is terminal. Successful flashlight exposure repels Gudo; interrupted exposure decays rather than instantly resetting.}
\filecard{SecretGame}{\href{https://github.com/itsnold/FiveNightsAtGudos/blob/e9f0ff45350572e7ba5fd0ff6185ab1251514e2f/src/main/java/com/gudos/game/minigame/SecretGame.java\#L31}{\nolinkurl{game/minigame/SecretGame.java}:31}}{Runs the hidden platformer rules without owning Swing or audio.\par{\sffamily\scriptsize\color{Muted}\textbf{Find:} action, jump, update, bundleHidden/witnessBeam}}{SecretPanel calls jump/action/update and draws read-only model values. The model integrates motion, patrols, bundle visibility and the pickup-triggered timer.}{Twenty-millisecond substeps reduce missed collisions. Witnesses detect the bundle center, not the player body. The same beam polygon is used for drawing and detection.}
\filecard{SecretPanel}{\href{https://github.com/itsnold/FiveNightsAtGudos/blob/e9f0ff45350572e7ba5fd0ff6185ab1251514e2f/src/main/java/com/gudos/ui/SecretPanel.java\#L35}{\nolinkurl{ui/SecretPanel.java}:35}}{Controls and renders the scrolling pixel-art minigame.\par{\sffamily\scriptsize\color{Muted}\textbf{Find:} bindings, tick, pause, silenceOnOutcome, draw helpers}}{GamePanel paper trigger -> Application.openSecret -> SecretPanel; key events call SecretGame; tick advances it; outcome freezes, silences and reports after three seconds.}{Use an offscreen 640x360 image enlarged nearest-neighbor. World scrolling is isolated from fixed HUD. reported and outcomeSilenced prevent duplicate callbacks/audio.}

\clearpage\pdfbookmark[0]{Sixty practice questions}{group-p:qa-1}\phantomsection\label{p:qa-1}\pdfbookmark[1]{Practice questions: 1--10}{bm-p:qa-1}\section*{Practice questions: 1--10}
\q{1. What is the main objective?}{Survive a six-minute active night until 6 AM while conserving power and defending against Gudo. {\scriptsize\color{Muted}\textit{Locate:} GameSession.update.}}
\q{2. Why use Java Swing?}{The implementation uses Java's desktop UI and drawing APIs to keep the game a self-contained offline desktop application. {\scriptsize\color{Muted}\textit{Locate:} Main.launch / ScenePanel.}}
\q{3. Where does execution begin?}{com.gudos.Main.main schedules launch on the Swing EDT. {\scriptsize\color{Muted}\textit{Locate:} Main.main / SwingWorker.}}
\q{4. Is GamePanel the entire game logic?}{No. It translates input and draws; GameSession and GudoEncounter enforce normal-night rules. {\scriptsize\color{Muted}\textit{Locate:} GamePanel and GameSession.}}
\q{5. Explain New Game as a facade.}{The button calls one Navigation.startNight(1) operation. Application hides the checks, new session/view creation, callback wiring, navigation and audio coordination. {\scriptsize\color{Muted}\textit{Locate:} MenuPanel -> Application.startNight.}}
\q{6. Does the interface itself make it a facade?}{The interface gives the small contract; the implementation's hidden subsystem coordination is the evidence of the facade. {\scriptsize\color{Muted}\textit{Locate:} Application.startNight.}}
\q{7. Why no separate newGame method?}{New Game is a label. Starting Night 1 is the same operation as starting a selected eligible night with argument one. {\scriptsize\color{Muted}\textit{Locate:} MenuPanel / NightSelectPanel / ResultPanel.}}
\q{8. Does New Game delete progress?}{It starts a new Night 1 attempt and preserves records/unlocks. {\scriptsize\color{Muted}\textit{Locate:} Application.startNight / Database.saveRun.}}
\q{9. What exactly does Continue continue?}{Progression, not a live checkpoint. It starts a fresh attempt of progress.nextNight. {\scriptsize\color{Muted}\textit{Locate:} MenuPanel constructor.}}
\q{10. What are your structural patterns?}{Facade at navigation/audio boundaries and a custom Decorator for CCTV effects. {\scriptsize\color{Muted}\textit{Locate:} Application / AudioService / rendering package.}}

\guidepage{Practice questions: 11--20}{p:qa-2}
\q{11. Explain the audio facade beyond its name.}{One playback request hides decoding, clip caching, volume/pan controls, looping, worker scheduling and cleanup. {\scriptsize\color{Muted}\textit{Locate:} AudioService.play / clip.}}
\q{12. Show the Decorator requirements.}{Every wrapper implements SceneRenderer, holds another SceneRenderer, delegates its render, then adds an effect. {\scriptsize\color{Muted}\textit{Locate:} RendererDecorator / render overrides.}}
\q{13. Why not just subclass a renderer for every effect combination?}{Composition lets us stack and order effects without creating a new subclass for each combination. {\scriptsize\color{Muted}\textit{Locate:} GamePanel constructor.}}
\q{14. Does decorator order matter?}{Yes. Dimming after noise darkens the noise; noise drawn after dimming appears over the darkened room. {\scriptsize\color{Muted}\textit{Locate:} GamePanel renderer chain.}}
\q{15. Why copy Graphics2D?}{An effect can change paint/composite/hints on its own context and dispose it without leaking that effect state into later HUD rendering. {\scriptsize\color{Muted}\textit{Locate:} VhsRendererDecorator.render.}}
\q{16. Why does reduced mode still render the room?}{VHS delegates first, then returns before adding optional noise. GamePanel also skips the switching chain. {\scriptsize\color{Muted}\textit{Locate:} GamePanel.renderCamera.}}
\q{17. Explain the Builder with real participants.}{SecretLevelBuilder collects named level parts; Application.createSecretLevel is the production recipe; build creates SecretLevel, which SecretGame uses. {\scriptsize\color{Muted}\textit{Locate:} Application.createSecretLevel.}}
\q{18. Is the level fully immutable?}{It copies construction inputs and lists, but Rectangle objects exposed through accessors remain mutable. {\scriptsize\color{Muted}\textit{Locate:} SecretLevel constructor.}}
\q{19. Is this the State pattern?}{It uses finite state-machine logic, but not separate polymorphic GoF State objects. {\scriptsize\color{Muted}\textit{Locate:} GudoEncounter / SecretGame.State.}}
\q{20. Where is encapsulation?}{GameSession fields are private; controls change them through guarded methods. {\scriptsize\color{Muted}\textit{Locate:} GameSession.setLight.}}

\guidepage{Practice questions: 21--30}{p:qa-3}
\q{21. Where is runtime polymorphism?}{Application can tick/render through a ScenePanel reference, and GamePanel can call a SceneRenderer whether it is plain or decorated. {\scriptsize\color{Muted}\textit{Locate:} ScenePanel / SceneRenderer.}}
\q{22. Are you using a Singleton?}{There is one shared AssetManager in the normal boot flow, but no enforced Singleton class. {\scriptsize\color{Muted}\textit{Locate:} AssetManager constructor.}}
\q{23. Why callbacks instead of giving a panel all services?}{The panel can request just the behavior it needs, such as sound or pause, without owning storage/audio/window lifecycles. {\scriptsize\color{Muted}\textit{Locate:} GamePanel constructor callbacks.}}
\q{24. How does a button actually start anything?}{GameButton installs an ActionListener that runs the supplied Runnable. The screen passes the actual navigation/model callback. {\scriptsize\color{Muted}\textit{Locate:} GameButton constructor.}}
\q{25. How is a night timed?}{Application supplies measured active deltas; GameSession accumulates them to its normal 360-second duration. {\scriptsize\color{Muted}\textit{Locate:} Application.tick / GameSession.}}
\q{26. Why not subtract one fixed amount per frame?}{Frame delivery can vary, so time and power use actual supplied seconds and exact event-boundary steps. {\scriptsize\color{Muted}\textit{Locate:} GameSession.update.}}
\q{27. What if attack and morning happen together?}{The model explicitly gives morning priority, then power failure, then attack. {\scriptsize\color{Muted}\textit{Locate:} GameSession.update condition order.}}
\q{28. How does power work?}{0.085 percentage points per active second, plus 0.22 for the monitor or 0.20 for light. {\scriptsize\color{Muted}\textit{Locate:} GameSession.drainPerSecond.}}
\q{29. Is Gudo random or AI-trained?}{He follows deterministic route and timing rules, not a trained model. {\scriptsize\color{Muted}\textit{Locate:} GudoEncounter.advance.}}
\q{30. Why does light fail while I am holding F?}{The monitor must be down, the session active and the office view facing the current approach. Then exposure must reach its threshold before the attack deadline. {\scriptsize\color{Muted}\textit{Locate:} GameSession.update / Approach.faces.}}

\guidepage{Practice questions: 31--40}{p:qa-4}
\q{31. Why does he reappear when I look away?}{Interrupted effective exposure decays at 0.65 per second, so the fade reverses toward the occupied room. {\scriptsize\color{Muted}\textit{Locate:} GudoEncounter.advance.}}
\q{32. Why can't he appear in two camera rooms at once?}{The encounter returns one current location; sceneState checks that location for occupancy. {\scriptsize\color{Muted}\textit{Locate:} GameSession.sceneState.}}
\q{33. How does CCTV move when I am not watching?}{Its pan is a pure function of camera identity and session elapsed time, not accumulated viewing time. {\scriptsize\color{Muted}\textit{Locate:} CameraMotion.pan.}}
\q{34. Is the office actually 3D?}{The visual effect projects a wide 2D panorama using angle-based sampling. There is no 3D model/mesh here. {\scriptsize\color{Muted}\textit{Locate:} PanoramaRenderer.prepare.}}
\q{35. Why separate clean and occupied renderers?}{Each PanoramaRenderer returns its reused mutable output. Two instances keep both projected images available during the blend. {\scriptsize\color{Muted}\textit{Locate:} GamePanel.renderOffice.}}
\q{36. How does the paper click remain aligned during camera sweep?}{We reverse the window-to-scene transform, then reverse the current projection to a source image pixel and test the source-fraction hitbox. {\scriptsize\color{Muted}\textit{Locate:} GamePanel.mouseClicked.}}
\q{37. What starts the hidden timer?}{The first successful bundle pickup, not opening the minigame. {\scriptsize\color{Muted}\textit{Locate:} SecretGame.action.}}
\q{38. What exactly do witnesses detect?}{The bundle center, carried or dropped, inside their direction-facing triangle unless a cover rectangle contains that center. {\scriptsize\color{Muted}\textit{Locate:} SecretGame.update / witnessBeam.}}
\q{39. How do you win the minigame?}{Carry the bundle into the stockroom destination and press E. Delivery is checked before dropping. {\scriptsize\color{Muted}\textit{Locate:} SecretGame.action / finishSecret.}}
\q{40. Does a secret win unlock the next night?}{No. Only normal Outcome.COMPLETED executes the unlock update. {\scriptsize\color{Muted}\textit{Locate:} Database.saveRun.}}

\guidepage{Practice questions: 41--50}{p:qa-5}
\q{41. How does minigame collision handle a long frame?}{SecretGame splits its supplied delta into steps no larger than0.02 seconds to reduce missed platforms or beam crossings. {\scriptsize\color{Muted}\textit{Locate:} SecretGame.update.}}
\q{42. Why SQLite and also text files?}{SQLite is authoritative relational progress/history; settings are human-editable preferences; the log is a readable report. {\scriptsize\color{Muted}\textit{Locate:} AppPaths and store classes.}}
\q{43. Explain the database relations.}{One profile has one progress summary and many night\_runs. progress.profile\_id is its primary/foreign key; each run has its own UUID and a profile foreign key. {\scriptsize\color{Muted}\textit{Locate:} Database.runsTable / schema.}}
\q{44. Why PreparedStatement?}{Player values are bound separately from SQL syntax, including quoted names. {\scriptsize\color{Muted}\textit{Locate:} Database.profile / PreparedStatement.}}
\q{45. What is atomic about saving?}{The attempt insert and unlock update commit or roll back together inside SQLite. {\scriptsize\color{Muted}\textit{Locate:} Database.saveRun.}}
\q{46. Are SQLite and shift-log.txt one transaction?}{No. They are separate writes coordinated by a retryable application workflow. {\scriptsize\color{Muted}\textit{Locate:} Application.savePendingRun.}}
\q{47. How do you stop duplicate records?}{A stable run UUID is the night\_runs primary key with ON CONFLICT DO NOTHING; the text report checks the same UUID marker. {\scriptsize\color{Muted}\textit{Locate:} RunRecord.id / store conflict/marker.}}
\q{48. Why not advance progress on every saved attempt?}{An attempt count includes failures and secrets. Progression represents normal survival, so the update requires a newly inserted COMPLETED result. {\scriptsize\color{Muted}\textit{Locate:} Database.saveRun.}}
\q{49. Why is the table limited to 200?}{The latest 200 rows keep the records view manageable; older rows are retained and totals still count them. {\scriptsize\color{Muted}\textit{Locate:} Database.loadRuns / loadProgress.}}
\q{50. Are profile names secure accounts?}{No. They select local records on this computer. There is no authentication or online service. {\scriptsize\color{Muted}\textit{Locate:} Settings / Database.profile.}}

\guidepage{Practice questions: 51--60}{p:qa-6}
\q{51. What happens if settings values are invalid?}{Parsing uses defined fallbacks and Settings sanitizes/clamps the final values. IO failures still report errors. {\scriptsize\color{Muted}\textit{Locate:} Settings constructor / store.load.}}
\q{52. Why a temporary settings file?}{Write a complete replacement first, then swap it over the real file, atomically where supported. {\scriptsize\color{Muted}\textit{Locate:} SettingsStore.save.}}
\q{53. Why can't an unexpected window focus loss keep the flashlight on?}{Application pauses the active panel and clears held controls; GameSession.pause forces light off. {\scriptsize\color{Muted}\textit{Locate:} windowLostFocus / releaseControls.}}
\q{54. Does closing always save?}{Normal close/quit asks and waits to save an active abandoned attempt. Force-kill or a process crash cannot guarantee that save. {\scriptsize\color{Muted}\textit{Locate:} Application.quit / close.}}
\q{55. What if the audio device is missing?}{The failed cue is remembered and playback returns without stopping gameplay. {\scriptsize\color{Muted}\textit{Locate:} AudioService.clip.}}
\q{56. How are resources packaged for offline play?}{Maven Shade bundles dependencies and runtime resources in the distributable JAR. AssetManager opens images, fonts and audio through classpath resources, so the built game does not need live asset downloads. {\scriptsize\color{Muted}\textit{Locate:} pom.xml Shade configuration and AssetManager.resource.}}
\q{57. Where should a new camera image/state be added?}{Put the image under the resources image folder, add a scene.state mapping in scenes.properties and use the state through the model. {\scriptsize\color{Muted}\textit{Locate:} scenes.properties / AssetManager.}}
\q{58. How would you add another enemy?}{It would require extending the model for multiple independent encounters/occupancy and then adapting rendering, cues and state updates. {\scriptsize\color{Muted}\textit{Locate:} GameSession encounter ownership.}}
\q{59. What are the main limitations?}{Local-only profiles, no mid-shift restore, one deterministic enemy, hardware-dependent audio, basic minigame input validation, and separate DB/text writes. {\scriptsize\color{Muted}\textit{Locate:} Limits page.}}
\q{60. How do you know the added comments did not change the game?}{The cleaned comments retain identical parsed executable syntax, and all 70 production class files still match the original byte-for-byte with debug metadata disabled. {\scriptsize\color{Muted}\textit{Locate:} Validation note on limits page.}}

\clearpage\pdfbookmark[0]{Rehearsal and reference}{group-p:rehearse}\phantomsection\label{p:rehearse}\pdfbookmark[1]{A five-minute rehearsal and common snags}{bm-p:rehearse}\section*{A five-minute rehearsal and common snags}

\par\noindent\begin{tabularx}{\linewidth}{L{.13\linewidth}L{.34\linewidth}X}\toprule
\textbf{Time}&\textbf{Do}&\textbf{Explain}\\\midrule
0:00--0:30&Show title/project folders&Offline Swing game; separate model, screen and coordinator.\\
0:30--1:00&Click New Game; open startNight&One navigation facade hides checks and model/view/audio coordination. History stays.\\
1:00--1:40&Pan, monitor, camera, light&Model rejects incompatible devices; correct facing/exposure matters. CCTV uses session time.\\
1:40--2:20&Pause; show update/encounter&Active-time freeze and exact event deadlines. Morning wins an exact tie.\\
2:20--3:00&Open renderer construction&Same SceneRenderer contract supports nested effects; draw order matters.\\
3:00--3:45&Confirm abandonment; show Records&One final UUID receipt; background saves; newest 200 versus total.\\
3:45--4:20&Show saveRun and log marker&SQLite transaction is separate from report append; ordinary retry avoids duplicates.\\
4:20--5:00&Show level recipe / hidden path&Builder product; pickup-triggered timer; both secret outcomes restart same normal night.\\\bottomrule\end{tabularx}\par
\subsection*{If something does not match expectations}
\begin{itemize}
\item \textbf{Continue disabled}: this profile has no recorded attempt. \textbf{Night locked}: check name and resolved data folder.
\item \textbf{F ineffective}: monitor down, active session, face current doorway, enough light before deadline. Wrong direction still drains.
\item \textbf{Frequent pause}: the window is losing focus. \textbf{Reduced camera darkness}: SIGNAL LOST still hides a disrupted feed for two active seconds.
\item \textbf{Secret trigger fails}: CAM04, actual paper, five valid left-clicks, Nights 1--4. \textbf{Timer keeps running in cover}: expected after first pickup.
\item \textbf{Save failed but row exists}: database may have committed before text; Retry keeps the same UUID. \textbf{Fewer table rows than totals}: newest 200 limit, not deletion.
\item \textbf{No sound}: check volume/mixer and stderr. \textbf{Build download failure}: dependency access is needed once; the built JAR itself is offline.
\end{itemize}
\note{Prepare the presentation machine}{Build before the session, use an isolated demo folder, pre-open key files, verify actual speakers/focus/fullscreen, and decide whether demonstrating a six-minute win or an honest abandonment/loss.}

\guidepage{Implementation limits and source evidence}{p:limits}

\subsection*{Boundaries to explain accurately}
\begin{enumerate}
\item \textbf{No mid-shift restore.} Saved runs are ended receipts, not complete live-session serialization. Continue starts fresh.
\item \textbf{Separate stores.} Attempt/unlock is transactional within SQLite; text append is separate. UUID retries address normal partial success, not all crash/journaling cases.
\item \textbf{Local profiles.} Names are labels on one computer, without passwords, online identity or synchronization. File access permits external local edits.
\item \textbf{One deterministic enemy.} Multiple independent enemies need model/occupancy/render/cue changes; they are not already supported.
\item \textbf{Input/config assumptions.} SecretGame assumes valid caller deltas; Builder checks selected constraints. Returned level Rectangles remain mutable despite input copies.
\item \textbf{Time and hardware.} UI delta cap can slow simulation during stalls. Optional mixer controls vary; silent fallback is deliberate.
\end{enumerate}
\subsection*{What the verification establishes}
All 45 application Java files parse with unchanged executable syntax after comment cleanup. All 70 production class files match original byte-for-byte when compiled with debug metadata disabled, using the available JDK 21 compiler with source/target 17. Whitespace checks pass. This is not a full --release 17 API check or a run on an installed JDK 17 runtime.
Maven packaging could not download dependencies because repo.maven.apache.org was unreachable. No complete Maven package or native-window/speaker verification is claimed. The document uses actual headless Swing renders with controlled state. The PDF is compiled from this edition's XeLaTeX source.
\subsection*{Version and authority}
Repository: \url{https://github.com/itsnold/FiveNightsAtGudos}. Executable base: \code{b27297ab9e687518fc1543908740e5926969cf96}. Application source: \code{e9f0ff45350572e7ba5fd0ff6185ab1251514e2f}. Links pin to that immutable source revision; later documentation-only commits do not alter those line references.
Current Java, active scene manifest, Sound enum, pom/launcher configuration determines actual behavior. Some older README/asset names retain obsolete dining or synthesized-audio descriptions. Benefits/decisions in this guide are reasoned from code, not undocumented author interviews.
\note{Where the exhaustive detail went}{Every explicitly declared method/constructor still has its explanation beside the code. This edition keeps all 45 application Java-file responsibilities and the important algorithms, decisions, navigation and questions while avoiding repetitive per-getter narration in the PDF.}

\guidepage{Alphabetical Java file locator}{p:index}
{\small Click a page number to open its file card. All 45 application files live under \code{src/main/java/com/gudos}.}\par\vspace{3pt}
\begin{minipage}[t]{.485\linewidth}\small\par\noindent\begin{tabularx}{\linewidth}{Xr}\toprule\textbf{Java file}&\textbf{Page}\\\midrule
AppPaths.java&\hyperref[file:AppPaths]{\pageref*{file:AppPaths}}\\
Application.java&\hyperref[file:Application]{\pageref*{file:Application}}\\
AssetManager.java&\hyperref[file:AssetManager]{\pageref*{file:AssetManager}}\\
AudioService.java&\hyperref[file:AudioService]{\pageref*{file:AudioService}}\\
BackdropPanel.java&\hyperref[file:BackdropPanel]{\pageref*{file:BackdropPanel}}\\
CameraMotion.java&\hyperref[file:CameraMotion]{\pageref*{file:CameraMotion}}\\
CameraSceneRenderer.java&\hyperref[file:CameraSceneRenderer]{\pageref*{file:CameraSceneRenderer}}\\
Chiptune.java&\hyperref[file:Chiptune]{\pageref*{file:Chiptune}}\\
Database.java&\hyperref[file:Database]{\pageref*{file:Database}}\\
DimmedRendererDecorator.java&\hyperref[file:DimmedRendererDecorator]{\pageref*{file:DimmedRendererDecorator}}\\
ExtrasPanel.java&\hyperref[file:ExtrasPanel]{\pageref*{file:ExtrasPanel}}\\
GameButton.java&\hyperref[file:GameButton]{\pageref*{file:GameButton}}\\
GamePanel.java&\hyperref[file:GamePanel]{\pageref*{file:GamePanel}}\\
GameSession.java&\hyperref[file:GameSession]{\pageref*{file:GameSession}}\\
GudoEncounter.java&\hyperref[file:GudoEncounter]{\pageref*{file:GudoEncounter}}\\
InfoPanel.java&\hyperref[file:InfoPanel]{\pageref*{file:InfoPanel}}\\
JumpscarePanel.java&\hyperref[file:JumpscarePanel]{\pageref*{file:JumpscarePanel}}\\
LoadingPanel.java&\hyperref[file:LoadingPanel]{\pageref*{file:LoadingPanel}}\\
Main.java&\hyperref[file:Main]{\pageref*{file:Main}}\\
MenuPanel.java&\hyperref[file:MenuPanel]{\pageref*{file:MenuPanel}}\\
Navigation.java&\hyperref[file:Navigation]{\pageref*{file:Navigation}}\\
NightSelectPanel.java&\hyperref[file:NightSelectPanel]{\pageref*{file:NightSelectPanel}}\\
Outcome.java&\hyperref[file:Outcome]{\pageref*{file:Outcome}}\\
\bottomrule\end{tabularx}\par\end{minipage}\hfill\begin{minipage}[t]{.485\linewidth}\small\par\noindent\begin{tabularx}{\linewidth}{Xr}\toprule\textbf{Java file}&\textbf{Page}\\\midrule
PanoramaRenderer.java&\hyperref[file:PanoramaRenderer]{\pageref*{file:PanoramaRenderer}}\\
PausePanel.java&\hyperref[file:PausePanel]{\pageref*{file:PausePanel}}\\
Progress.java&\hyperref[file:Progress]{\pageref*{file:Progress}}\\
RecordsPanel.java&\hyperref[file:RecordsPanel]{\pageref*{file:RecordsPanel}}\\
RendererDecorator.java&\hyperref[file:RendererDecorator]{\pageref*{file:RendererDecorator}}\\
ResultPanel.java&\hyperref[file:ResultPanel]{\pageref*{file:ResultPanel}}\\
RunRecord.java&\hyperref[file:RunRecord]{\pageref*{file:RunRecord}}\\
Scene.java&\hyperref[file:Scene]{\pageref*{file:Scene}}\\
ScenePanel.java&\hyperref[file:ScenePanel]{\pageref*{file:ScenePanel}}\\
SceneRenderer.java&\hyperref[file:SceneRenderer]{\pageref*{file:SceneRenderer}}\\
SecretGame.java&\hyperref[file:SecretGame]{\pageref*{file:SecretGame}}\\
SecretLevel.java&\hyperref[file:SecretLevel]{\pageref*{file:SecretLevel}}\\
SecretLevelBuilder.java&\hyperref[file:SecretLevelBuilder]{\pageref*{file:SecretLevelBuilder}}\\
SecretPanel.java&\hyperref[file:SecretPanel]{\pageref*{file:SecretPanel}}\\
Settings.java&\hyperref[file:Settings]{\pageref*{file:Settings}}\\
SettingsPanel.java&\hyperref[file:SettingsPanel]{\pageref*{file:SettingsPanel}}\\
SettingsStore.java&\hyperref[file:SettingsStore]{\pageref*{file:SettingsStore}}\\
ShiftLogStore.java&\hyperref[file:ShiftLogStore]{\pageref*{file:ShiftLogStore}}\\
Sound.java&\hyperref[file:Sound]{\pageref*{file:Sound}}\\
Theme.java&\hyperref[file:Theme]{\pageref*{file:Theme}}\\
VhsOverlay.java&\hyperref[file:VhsOverlay]{\pageref*{file:VhsOverlay}}\\
VhsRendererDecorator.java&\hyperref[file:VhsRendererDecorator]{\pageref*{file:VhsRendererDecorator}}\\
\bottomrule\end{tabularx}\par\end{minipage}
\end{document}
